\documentclass[lettersize,journal]{IEEEtran}
\usepackage{amsmath,amsfonts}
\usepackage{algorithm}
\usepackage{algpseudocode}
\usepackage{array}
\usepackage[caption=false,font=normalsize,labelfont=sf,textfont=sf]{subfig}
\usepackage{textcomp}
\usepackage{stfloats}
\usepackage{url}
\usepackage{graphicx}
\usepackage{cite}
\usepackage{multirow}
\usepackage{amssymb}
\usepackage{amsthm}
\usepackage{xcolor}
\usepackage{booktabs}
\usepackage{tikz}
\usetikzlibrary{positioning}
\usepackage{pgfplots}
\usepackage{needspace}
\pgfplotsset{compat=1.17}
\usepackage{flushend}

\newtheorem{theorem}{Theorem}
\newtheorem{proposition}[theorem]{Proposition}
\newtheorem{definition}{Definition}
\newtheorem{remark}{Remark}

\hyphenation{op-tical net-works semi-conduc-tor IEEE-Xplore}

\begin{document}

\title{ECLIPSE: Per-Record Membership Encoding Passes Behavioral\\
 Unlearning Audits}

\author{Weimin~Lai,~Qiao~Yan,~and~Sizhe~Liang%
\thanks{Manuscript draft, 2026. The authors are with the College of
Computer Science and Software Engineering, Shenzhen University,
Shenzhen, China.}}

\markboth{IEEE Transactions on Dependable and Secure Computing}%
{IEEE Transactions on Dependable and Secure Computing}

\maketitle

\begin{abstract}
Machine unlearning answers a data-deletion request by
returning a model that no longer reflects the deleted
records, and behavioral deletion audits certify compliance by
comparing aggregate statistics of the returned model against
honest retraining within tolerance bands. Existing attacks
bypass other verification layers or operate without an
acceptance constraint; none must satisfy a live, calibrated,
multi-metric acceptance rule while the encoding takes place.
This paper shows that the rule itself can be turned into a
carrier: we propose ECLIPSE, a provider-side construction
whose feasible-direction contrastive encoder optimizes under
the audit's acceptance test. In the principal
fine-tuning-referenced setting, a
provider with access to the realized audit scorer constructs
accepted models with increased membership-ranking signal;
stronger working-point readout requires an honest reference
ensemble. We further trace why the audit tolerates the
encoding: the acceptance region is centered on the reference
operator's own forget-fit, the utility-side allowance is the
resource the encoder spends, and the tested constructions do
not reproduce the signal under from-scratch-referenced audits.
Four audit-side countermeasures are evaluated; adaptive
encoders survive the joint covariance rule and threshold
tightening, while the aligned statistic and the aggregate
z-detector hold. Transfer and persistence tests bound the
effect across reference families, classes, deletion scales,
architectures, and benign fine-tuning.
\end{abstract}

\begin{IEEEkeywords}
machine unlearning, deletion audits, membership inference,
conformal calibration, adversarial providers
\end{IEEEkeywords}


\section{Introduction}\label{sec:intro}
% ============================================================
% Para 1: Domain background → research problem + core challenge
% Check 1 fix: metric list simplified to 3 representative items
% ============================================================
\IEEEPARstart{M}{achine} unlearning has emerged as the primary
mechanism for honoring data-deletion requests in machine learning
services: instead of retraining from scratch, the provider returns
a model claimed to no longer reflect the deleted records. Because
the requester cannot inspect the training pipeline, certification
rests on \emph{behavioral deletion audits}. The auditor queries
the returned model, computes a suite of aggregate statistics
(accuracy, forget-set loss, membership-inference advantage, and
related metrics), and accepts the model iff every statistic falls
within a tolerance band calibrated from honest retraining
\cite{goel2022adversarial,tu2025reliable,dorna2025openunlearning,xue2025survey}.
The audit is the \emph{only} barrier between a dishonest provider
and the data subject's erasure right: the provider must return
some model, cannot be compelled to retrain honestly, and no one
outside the service observes the training run. This paper
addresses the following research questions:

\begin{itemize}
\item[\textbf{RQ1}] \emph{Feasibility:} Can a dishonest
provider construct a model that passes every audit statistic
while carrying readable per-record membership signals?
\item[\textbf{RQ2}] \emph{Mechanism:} Why does the audit
tolerate such encoding---is the tolerance structural or an
implementation artifact?
\item[\textbf{RQ3}] \emph{Defense:} Can audit-side hardening
detect or block the encoding, and at what cost?
\end{itemize}

% ============================================================
% Para 2: Existing research directions → their shortcomings
% ============================================================
A recent survey names the open threat class: adversaries who
``manipulate model behavior to appear compliant during audit,
while covertly retaining information''~\cite{xue2025survey}.
Existing attacks either bypass a different verification
interface or operate without any acceptance constraint: the
dishonest-provider study of Zhang et
al.~\cite{zhang2024fragile} circumvents replay-based and
backdoor-based protocols, and information encoding by
malicious trainers---Song et al.~\cite{song2017remember},
Truth Serum~\cite{truthserum2022}, Privacy
Backdoors~\cite{wen2024privacy}---passes at most a
model-quality gate, never a live, calibrated, multi-metric
acceptance rule \emph{while} the encoding takes place. The
unexamined surface is the acceptance region's
\emph{interior}: what information a dishonest provider can
deliberately place there, and who can read it after acceptance.

% ============================================================
% Para 3: Our method → how it solves the above problems
% Check 2 fix: no specific numbers; qualitative statement only
% ============================================================
This paper proposes ECLIPSE (\emph{Encoding that Crosses audIt
Lines while Preserving Signals of Erasure-eligible records}), a
provider-side construction that answers RQ1 affirmatively.
ECLIPSE exploits a structural property of the
acceptance rule: the accuracy component leaves an allowance that
the provider can spend, while the deletion itself forces the
forget-side tolerance wide---room the encoder, as we show,
does not use (\S\ref{sec:allowance}). Implementing the
encoding faces two challenges. First, acceptance is live:
starting from a feasible model, FIDE commits only candidates
that pass the unchanged scorer for every audited statistic
simultaneously; trial candidates may be rejected. Our
feasible-direction contrastive encoder puts the audit
check inside the optimization loop, so every committed
iterate is feasible and the encoder may be terminated at any
point. Second, membership must be readable \emph{per record}
after acceptance, not merely detectable in aggregate. Our
reference-calibrated readout (CAL-Z) calibrates each record
against the audit's own honest references and reads it at a
working-point threshold once the model is accepted and
deployed (Section~\ref{sec:method}).

% ============================================================
% Para 4: Main contributions (bullets) with key numbers
% (numbers live here, not in P3 — TDSC convention)
% Check 5 fix: "false alarm rates" → "false-rejection rates"
% ============================================================
Our contributions are as follows.
\begin{itemize}
\item \emph{ECLIPSE and the acceptance--readout gap (RQ1).}
The first construction of deliberate per-record membership
encoding inside a live, conformally calibrated multi-metric
deletion audit. Nine of nine starts show label-blind AUC
gains (all sensitivity intervals positive); a
reference-calibrated readout reaches per-record TPR $0.43$
at a nominal $5\%$ threshold---from a model that passes
every audit statistic.
\item \emph{The two-part mechanism (RQ2).} \emph{Admissibility:}
under a retrain-calibrated audit both tested F-aware starts are
rejected---the original model on the deletion statistics
(forget-set loss $0.10$ against a centre of $0.82$) while its
accuracy is in band, the fine-tuned start on accuracy and
margin---and the F-free starts it admits deliver no signal, so
the retrain boundary closes at admission as well as through
the small allowance; because reference family and start
ancestry are collinear, the two closings are reported jointly,
not as one identified cause. \emph{Budget:} within the admitted
family, delivered gain rises monotonically with the utility
allowance ($0.084$ to $0.253$ over a six-fold range; the two
tightest settings are single-start probes) and is invariant to
the forget-side width---the width deletion-inflated dispersion
widens (\S\ref{sec:allowance}). The allowance is not
sufficient; class and scale remain bounded empirically.
\item \emph{Measurement discipline for deletion-audit
readings.} Matched-budget selection and logit-scale controls
that reinterpret recovery scores, rollback exclusion on a
full grid, and split-conformal coverage calibration with a
declared false-rejection target---tools that prevent both
inflated attack claims and false-rejection rates.
\item \emph{Defense evaluation with adaptive counter-arms
(RQ3).} The aligned seventh statistic
($+0.14\to+0.04$), the covariance-calibrated joint rule,
and the aggregate z-detector with held-out validation,
plus a stronger membership probe as an instrument-level
test. Outcomes differ by instrument and are not
interchangeable: the joint rule and threshold tightening are
each evaded by re-optimized encoders; the aligned statistic's
compression survives re-optimization and two objective
redesigns; the composite z-detector arm admits feasible finals
with no positive gain; and the stronger probe rejects the
frozen finals while a fresh encoding under it delivers again.
\end{itemize}

We evaluate ECLIPSE on CIFAR-10/100 with ResNet-18 and
MobileNetV2 under preregistered gates; the results
demonstrate a passing-yet-readable encoding whose tolerance
traces to the audit's calibration structure, with defense
outcomes that differ by instrument.

% ============================================================
% Para 5: Paper organization
% Check 4 fix: RQ-to-section mapping
% ============================================================
The remainder of this paper is organized as follows.
Section~\ref{sec:related} reviews related work.
Section~\ref{sec:threat} states the threat model and the
audit protocol. Section~\ref{sec:method} presents ECLIPSE
(RQ1: construction), together with the structural property
it exploits---deletion-forced dispersion of the acceptance
region. Section~\ref{sec:eval} evaluates ECLIPSE in three
acts: feasibility and delivered signal (RQ1), why the audit
tolerates the encoding (RQ2: mechanism), and the
auditor's countermeasures (RQ3). Section~\ref{sec:discussion}
states limitations and directions; Section~\ref{sec:conclusion}
concludes.

\section{Background and Related Work}\label{sec:related}
\subsection{Statistical deletion audits}
Auditing unlearning spans parametric replay (proof-of-unlearning
logs \cite{thudi2022necessity}), cryptographic and proof-carrying
variants \cite{xue2025survey}, and behavioral statistics. On the
behavioral side, Goel et al.~\cite{goel2022adversarial} introduced
the Interclass Confusion test as a stronger black-box check and
showed that poisoned labels survive naive unlearning evaluations;
Tu et al.~\cite{tu2025reliable} formalize MIA-based evaluation as a
cryptographic game, expose a split-sampling sabotage construction
against naive two-split testing, and derive a swap-based fix;
benchmark suites \cite{dorna2025openunlearning} document that
unlearning metrics disagree with one another, and the verification
survey \cite{xue2025survey} systematizes the field. Our audit
model---five deployable metrics with tolerance bands---is the
deployment-facing representative of this line, and we inherit its
central, largely unexamined design parameter: the tolerance. None of
these works characterizes the acceptance region that the tolerance
induces---nor, at record granularity, the information a dishonest
provider can park inside it while passing.

\subsection{Attacks on unlearning and its verification}
The known attack families occupy several layers: request-layer
degradation \cite{huang2025unlearn}, data-layer payloads
\cite{di2023hidden}, and log-layer forgery
\cite{thudi2022necessity}. Closest to us, Zhang et
al.~\cite{zhang2024fragile} study the dishonest provider
directly: adversarial unlearning procedures retain the
forgotten data while passing \emph{replay-based and
backdoor-based} verification protocols (their evaluation also
includes membership-inference measurements); their closest
analogues under a banded audit are measured in
Sections~\ref{sec:wash} and~\ref{sec:region}.
\emph{Active information encoding} by malicious trainers is
established in the broader security literature---Song et
al.~\cite{song2017remember} encode retrievable information
behind a model-quality gate, Truth
Serum~\cite{truthserum2022} amplifies membership inference
via poisoning, Privacy Backdoors~\cite{wen2024privacy} leak
membership through later fine-tuning---all before or during
training, without a live, deletion-referenced, calibrated
multi-metric acceptance constraint.

Our readout draws on a different lineage: the
membership-inference foundations
\cite{shokri2017membership, yeom2018privacy, salem2019mlleaks},
difficulty calibration \cite{watson2022importance}, and
reference-normalized decoding \cite{carlini2022membership};
Hayes et al.~\cite{hayes2024inexact} show that
population-level unlearning MIAs understate per-example
leakage. The present study differs in direction: a provider
\emph{intentionally constructs} an artifact that satisfies the
acceptance region while carrying per-record signal, working
\emph{under} a live conformal audit.
\emph{Post-hoc transformation} attacks exploit the
audit-deployment gap instead: quantization after unlearning
restores forgotten content \cite{zhang2025catastrophic}
without passing the audit's own statistics. The survey
\cite{xue2025survey} names behavioral-space forgery as a
threat class; on the defense side, Di et
al.~\cite{di2024adversarial} make unlearning robust to MIA
auditors---our results suggest the harder problem is not
making honest models pass, but making dishonest models fail.
Table~\ref{tab:neighbors} consolidates the delimitation; the
qualifier-scoped first claim of Section~\ref{sec:intro} is
made against these families and the unlearning-verification
and malicious-encoding search record documented in the
project archive (retrieval through 2026-11). Throughout, a
returned model that is optimized to pass the audit while
retaining forget-set influence is called a \emph{wash}.

\begin{table*}[t]
\caption{Attack families on unlearning and its verification, and
where this work sits. ``Passes audit statistics'' asks whether the
returned artifact must itself satisfy the audit's banded metrics
(forged logs and deployment-time transforms evade this constraint by
construction; the wash must meet it head-on).}
\label{tab:neighbors}
\centering
\small
\begin{tabular}{p{24mm}p{28mm}p{24mm}p{16mm}p{42mm}}
\toprule
Family & Representative & Artifact attacked & Passes audit stats? & What is characterized \\
\midrule
Request-layer & Unlearn-and-Burn \cite{huang2025unlearn} & the deletion request itself & --- & utility destruction via adversarial requests \\
Data-layer & Hidden Poison \cite{di2023hidden} & training data & --- & payload activated by the unlearning event \\
Log-layer & auditable definitions \cite{thudi2022necessity} & replayed training trajectories & no (forged logs) & process-verification forgery \\
Verification-bypass & fragile verification \cite{zhang2024fragile} & replay/backdoor verification protocols & yes (those protocols) & dishonest provider retains data past replay and backdoor checks \\
Statistical-test sabotage & split-sampling construction \cite{tu2025reliable} & the audit's evaluation split & yes (the sabotaged test) & evaluator-side sabotage of a naive two-split MIA evaluation (swap-based fix derived); aggregate sabotage, no per-record encoding, no calibrated multi-metric band \\
Post-hoc transform & quantization recovery \cite{zhang2025catastrophic} & deployment-time quantization & no (audit vs.\ deployment gap) & forgotten-content restoration at deployment \\
Behavioral forgery & fingerprint suppression, signal obfuscation, performance spoofing \cite{xue2025survey} & probe responses, privacy signals, metrics & partially & named threat class; mechanism-level open \\
\midrule
\textbf{In-band encoding (this work)} & first-order contrastive encoding under a live conformal audit & the returned model vs.\ banded statistical acceptance & \textbf{yes, by construction} & the acceptance--readout gap: per-record membership readable in passing models; forced tolerance; partial suppression \\
\bottomrule
\end{tabular}
\end{table*}

\subsection{Unlearning operators}
Exact retraining \cite{bourtoule2021machine}, certified
removal \cite{guo2020certified}, and approximate operators
\cite{foster2024ssd, golatkar2020eternal, kurmanji2023towards}
define what the honest provider would return; LLM-scale
variants are benchmarked in \cite{maini2024tofu,
dorna2025openunlearning}. Our threat model is orthogonal to
the operator: the dishonest provider returns \emph{whatever
passes}, so the audit's geometry is the object of study.
Across all families of Table~\ref{tab:neighbors}, the
difference from this work is threefold, as sharpened in
Section~\ref{sec:threat}. Firstly, \emph{live acceptance}:
the returned artifact must itself satisfy every banded
statistic, where forged logs and deployment-time transforms
evade this constraint by construction. Secondly,
\emph{provider-known calibration}: the provider holds the
realized scoring function, a stronger information setting
than the listed families assume. Thirdly, \emph{no process
control}: nothing about training is observed or attested, so
the attack surface is the returned model's behavior alone.


\section{Threat Model and Audit Protocol}\label{sec:threat}\label{sec:audit}
This section sets the stage in three steps: the parties,
capabilities, and definitions; the audit as a computational
object (metric suite, acceptance region, calibration); and the
protocol versions under which results are reported.

\subsection{Parties, Capabilities, and Definitions}
A \emph{provider} holds model $\theta_0$ trained on dataset $D$ and
receives a deletion request for a set $F \subset D$. An \emph{auditor}
receives the returned model $\theta'$, knows $F$ and the retained set
$R = D \setminus F$, and runs a statistical audit. The provider is
dishonest: it fully controls $\theta'$ and may retrain internally as
many times as it likes before returning one model. The auditor is
honest and black-box-plus-losses---predictions and per-sample losses
are all that an outside party can query---and retrains only to build
its reference ensemble.
Table~\ref{tab:capabilities} fixes the full capability split.

\begin{table}[t]
\caption{Capability split. The provider is strong (white-box over its
own returned parameters, audit-aware); the auditor is deployable
(black-box predictions and losses; no retraining of its own).}
\label{tab:capabilities}
\centering
\small
\begin{tabular}{p{40mm}p{19mm}p{15mm}}
\toprule
Capability & Provider & Auditor \\
\midrule
Returned parameters $\theta'$ & full control & metrics only \\
Audit suite and band semantics & known & defines \\
Internal retraining budget & $M$ retrains & --- \\
$F$, $R$ membership & knows & knows \\
Held-out queries / losses & has & has \\
Retraining compute & internal & reference-ensemble setup only \\
\bottomrule
\end{tabular}
\end{table}

\begin{definition}[Acceptance region]
\label{def:region}
Given metric center $c \in \mathbb{R}^5$ and band $b \in
\mathbb{R}^5_{>0}$, the acceptance region is
\begin{equation}
\label{eq:region}
\mathcal{A}(c,b) = \bigl\{\theta : |m_i(\theta) - c_i| \le b_i,\
i = 1,\dots,5\bigr\},
\end{equation}
where $c$ and $b$ are the median and dispersion radius (maximum
deviation from the median) of an honest retraining ensemble
$\{\theta_R^{(s)}\}_{s=1}^{S}$.
\end{definition}


\subsection{Adversary and Downstream Capability Tiers}
Every downstream claim is scoped to one of three capability
levels, stated here once. The \textbf{provider} (all tiers)
knows $\theta_0$, the training data, and $F$; may use any
retained records; is granted full audit-scoring capability
(centers, scales, threshold $\hat q_{0.20}$) with unlimited
offline feasibility queries; and submits one final model.
\textbf{Tier 1 --- label-blind inferrer} (deployable setting):
obtains the returned model, sees candidate records with class
labels but no membership labels, and fixes scoring direction
and threshold on auxiliary never-trained records. Under this
tier the demonstrated harm is \emph{ranking-level} membership
inference (AUC gains on 9/9 starts); working-point
identification is limited to TPR $\le 0.10$ at nominal $5\%$
FPR (Section~\ref{sec:working-points}).
\textbf{Tier 2 --- reference-equipped evaluator}: as Tier 1,
plus an honest reference ensemble (six models)---a party that
holds or can reconstruct the auditor's own reference
infrastructure. The reference-calibrated readout
(Section~\ref{sec:zreadout}, Table~\ref{tab:zreadout})
operates at this tier.
\textbf{Tier 3 --- auditor-side}: the auditor's own ensemble.
It is what makes the Tier-2 readout demonstrable, but it adds
no information to an auditor who already knows $F$ and $R$
(Section~\ref{sec:limitations}).


\subsection{Cost accounting for both sides}
The provider's dominant cost is internal reference generation:
the reliability figures in this paper assume $M = 6$ honest
processing runs of the retained set under the audit family's
operator (roughly six times the honest deletion-processing
cost), before an optimization that is gradient-cheap (at most
60 accepted steps on $|F|$ records per candidate). The
auditor's deployable budget is query-only. A
two-retrain pilot can measure $\hat\gamma$
(\S\ref{sec:region}) but has no prospective validation as a
prediction rule.



\subsection{Design Constraints for ECLIPSE}
ECLIPSE operates under three constraints that distinguish it from
every published attack family (Table~\ref{tab:neighbors}):
(1)~\emph{Live acceptance}: the returned model must satisfy every
banded metric simultaneously---the attack surface is the
\emph{interior} of the acceptance region.
(2)~\emph{Provider-known calibration}: the provider holds the
audit's centers, scales, and threshold ($\hat q_{0.20}$), and
may make unlimited offline feasibility queries.
(3)~\emph{No process control}: the provider controls only the
returned parameters $\theta'$.

\subsection{Audit suite}
We now define the audited object. Two labeling families appear
below: $A_5$--$A_8$ are audit-protocol versions
(Section~\ref{sec:protocol-versions}), and Tier 1--3 are the
downstream capability levels just stated.

Following deployable practice \cite{goel2022adversarial,
dorna2025openunlearning}, the audit computes five deployable metrics:
$m_1$ overall test accuracy; $m_2$ class accuracy on the deleted
records' class; $m_3$ mean cross-entropy on $F$; $m_4$ two-sample
Kolmogorov--Smirnov distance between per-sample losses on $F$ and a
frozen retained subsample (3{,}000 records, fixed across the study);
$m_5$ membership-inference AUC of a logistic probe on per-sample
(loss, correct-class-margin) features separating $F$ from the
retained subsample (random guessing $=$ 0.5). Metrics are
evaluated in the records' canonical
(unaugmented) form, matching the training input representation, so
that evaluation is a fixed forward pass on a given device
(boundary flips at the acceptance grid are device-disclosed
in Supplementary Table S12; the separate-execution base-AUC
difference of Table S6 remains unexplained). The model passes iff
$|m_i(\theta') - c_i| \le b_i$ for all $i$, with $(c,b)$ from the
honest ensemble (Definition~\ref{def:region}).
Section~\ref{sec:calibration} shows this ensemble calibration is
forced: bands calibrated on full-data retraining variance reject the
honest ensemble itself.


\subsection{The audit as a protocol}\label{sec:protocol}
The deployable audit we attack is the following interaction.
(1)~The auditor fixes the metric suite $m_{1..5}$, the forget
set $F$, and a retained probe sample $S$. (2)~It obtains an
honest reference ensemble $\{\theta_R^{(s)}\}_{s=1}^{S}$ by
retraining on $D \setminus F$ under the service's own recipe,
and publishes or holds the centers $c_i = \mathrm{med}_s\, m_i$
and bands $b_i = \max_s |m_i - c_i|$. (3)~The service returns
$\theta'$. (4)~The auditor evaluates $m_{1..5}(\theta')$ on
canonical-form records and accepts iff every deviation is
within band. Step (2) is where the dispersion measurement
bites (the ensemble must be re-estimated per deletion scale),
and step (4) is where the acceptance region is invoked.
Nothing in the protocol constrains how the service produces
$\theta'$---precisely the gap the attack occupies.

Two scope statements bound the protocol. First, the step-(2)
retrains are the audit's trusted setup, run by the auditor or
an attested third party. All experiments in this
paper use the \emph{provider-known calibration} model: the
provider is granted the exact fixed audit-scoring function,
including its realized calibration parameters (centers,
scales, and threshold), and may query it offline without
limit. This assumption is deliberately conservative: granting
the provider the realized calibration is the harder case for
the auditor. Common ensemble distributions alone would not
imply this: the independent-calibration probe of
Supplementary Table S22 shows the realized estimated and
true score systems differ. Under the \emph{independent
calibration} model, five preregistered estimation draws show
that acceptance survives broadly while signal delivery is
draw-dependent rather than systematically blocked (the full
counts are restated with the limitations;
Section~\ref{sec:limitations}, Supplementary Tables S22,
S28). Second, the audit certifies only the \emph{submitted and
deployed model's behavior}; a dishonest service can retain
copies elsewhere, and no model-side audit certifies their
absence.


\subsection{Audit protocol versions}\label{sec:protocol-versions}
Three acceptance rules are used in this paper, plus one joint
variant evaluated as a hardening step ($A_5$--$A_8$; a
version-by-experiment summary table is in the supplement,
Supplementary Table S42). All share the same base metrics
and honest-reference architecture; they differ in the number
of scored statistics, the calibration procedure, and the
knowledge granted to the provider. The $A_6$ audit is the
paper's principal object: the five metrics above plus $m_6$
(mean forget-set logit margin). Centers $c_i$ and scales
$\sigma_i$ are the median and MAD of the fit ensemble; the
joint score is
\begin{equation}
\label{eq:a6}
s_6(\theta) = \max_i \frac{|m_i(\theta) - c_i|}{\sigma_i},
\end{equation}
and the threshold $\hat q_\alpha$ is the
conformal quantile of the calibration ensemble's scores. The
$A_7$ audit adds
\begin{equation}
\label{eq:a7}
d(\theta) = \mathbb{E}_{N_a}\|\nabla_x \ell\| -
\mathbb{E}_F\|\nabla_x \ell\|
\end{equation}
(frozen on independent fit models) and recalibrates; computing $d$
requires per-record input gradients, extending the auditor
beyond black-box-plus-losses. The $m_5$ probe's complete
specification (features, standardization, training, and
in-sample evaluation) is fixed in Supplementary Table S45.


\section{ECLIPSE: Methodology}\label{sec:method}\label{sec:encoding}

\subsection{Notation}
Table~\ref{tab:notation} fixes the symbols used throughout.

\begin{table}[t]
\caption{Notation.}
\label{tab:notation}
\centering
\small
\begin{tabular}{cl}
\toprule
Symbol & Meaning \\
\midrule
$D$, $F$, $R$ & training set, forget (deleted) set, retained set \\
$\theta_0$, $\theta'$ & original model, returned (encoded) model \\
$m_i(\cdot)$ & audit metric $i$ ($i=1,\dots,6$) \\
$c_i$, $\sigma_i$ & reference center / scale for metric $i$ \\
$s(\theta)$ & joint audit score $= \max_i |m_i(\theta)-c_i|/\sigma_i$ \\
$\hat q_{\alpha}$ & conformal threshold at level $\alpha$ \\
$N_c$, $N_a$ & construction pool, seventh-statistic set \\
$N_{\rm aux}$, $N_{\rm eval}$ & threshold-calibration / evaluation sets \\
$\ell(\theta;x,y)$ & cross-entropy loss of $\theta$ on $(x,y)$ \\
$J_{\rm CE}(\theta)$ & contrastive objective (Eq.~\ref{eq:ce-objective}) \\
$z_i$ & calibrated readout score for record $i$ \\
$\mu_i$, $\sigma^z_i$ & per-record reference mean / std for $z_i$ \\
\bottomrule
\end{tabular}
\end{table}

\subsection{Problem Statement}
\label{sec:problem}
The provider's task is a constrained optimization problem over
model parameters, defined by the audit of
Section~\ref{sec:audit}:

\begin{definition}[Constrained encoding problem]
\label{def:encoding}
Given the audit's centers $c$, scales $\sigma$, and threshold
$\hat q_{\alpha}$, the forget set $F$, and a same-class
construction pool $N_c$, find
\begin{equation}
\label{eq:encoding}
\theta' = \arg\min_{\theta}\; J_{\rm CE}(\theta)
\quad\mathrm{s.t.}\quad s(\theta) \le \hat q_{\alpha},
\end{equation}
where $s$ is the joint audit score of Eq.~(\ref{eq:a6}) and
$J_{\rm CE}$ is the contrastive objective of
Eq.~(\ref{eq:ce-objective}).
\end{definition}

The feasible set is non-convex and non-smooth: $s$ is a maximum
of six standardized deviations, one of which ($m_1$, accuracy)
lives on a discrete grid, and the binding component can switch
along a trajectory. Three properties of the problem shape the
solution: (i) with positive scales, the max-score constraint
$s(\theta) \le \hat q$ is exactly the conjunction of marginal
bands $|m_i - c_i| \le \hat q\,\sigma_i$---its distinctive
feature is that the half-widths $\hat q\sigma_i$ are calibrated
\emph{jointly} through the conformal quantile, and every
component must be checked; (ii) the submitted model must
be feasible under the fixed scorer---infeasible candidates
require a repair stage before submission; (iii) the audit is
queryable offline without limit (Section~\ref{sec:threat}), so
the constraint can be checked inside the loop.

\begin{remark}[Anytime termination, a conditional invariant]
\label{rem:anytime}
Given a feasible initial model and an unchanged scoring
function, feasible-direction descent
(Section~\ref{sec:fide}) commits only passing candidates, so
every committed iterate $\theta^{(t)}$ remains feasible and the
provider may terminate at any point---under a compute cap,
on an interrupt, or after a target number of accepted
steps---and still return a model that passes. The invariant
is about feasibility only: trial candidates may fail the
check, and early termination guarantees neither signal
delivery nor a pinned endpoint. Probes that start outside a
new constraint (the tightened-threshold and
independent-calibration runs of Supplementary Table S22)
first restore feasibility through the same backtracking
acceptance test.
\end{remark}

\subsection{The Structural Resource: Deletion-Forced
Dispersion}\label{sec:region}\label{sec:calibration}
ECLIPSE is built on a structural property of the audit's
acceptance region, so we state it before the construction. A
naive audit calibrates tolerance on the variance of full-data
retraining. This fails, and the failure mode is mechanical.
While $F$ is in the training set, every retrain fits it, and
the cross-seed spread of forget-side statistics is the
ordinary noise of fitting easy, repeatedly-constrained
examples. Excluding $F$ removes that constraint: the optimizer
is now free to allocate capacity elsewhere, and how much
residual fit each seed leaves on the deleted records varies
much more. Dispersion on forget-side statistics is therefore
\emph{created by the deletion itself}---the very act the audit
is certifying inflates the tolerance it must grant.
Concretely: with bands from six full-data retrains, the
\emph{honest} unlearned ensemble is rejected on $3/3$, $2/3$,
and $1/3$ of its seeds ($k{=}30/300/1000$; three seeds per
scale); the full dispersion measurement and its
scale-by-metric detail are in Section~\ref{sec:eval}. The
ensemble band we use throughout is the minimal correction and
covers every honest model in the calibration ensemble by
construction; it is a sample statistic, not a coverage
guarantee---a deployed audit should fix a target
false-positive rate and verify coverage on fresh honest
outputs (the split-conformal analysis in
Section~\ref{sec:eval} addresses this directly).

\paragraph*{A linear-model mechanism.}
The property can be verified in closed form in the
linear-Gaussian setting. Let $y = X\beta^\ast + \varepsilon$ with $\varepsilon \sim
\mathcal{N}(0, \sigma^2 I_n)$, let $\hat\beta_{\rm full}$
minimize the squared loss on all $n$ rows and
$\hat\beta_{R}$ on the $n{-}k$ rows outside $F$, and let the
forget-side statistic be the mean squared residual
$m_F(\hat\beta) = \frac{1}{k}\|y_F - X_F\hat\beta\|^2$.

\begin{proposition}[Constraint release inflates forget-side
dispersion]
\label{prop:dispersion}
Over the noise draw $\varepsilon$ (all else fixed), let
$\lambda_F^{(R)} = \mathrm{tr}(X_F (X_R^\top
X_R)^{-1} X_F^\top)$ be the total leverage of the
deleted block against the retained design. Then,
whenever $\lambda_F^{(R)} > 0$,
\begin{equation}
\label{eq:disp}
\mathrm{Var}_\varepsilon\big[m_F(\hat\beta_{R})\big] -
\mathrm{Var}_\varepsilon\big[m_F(\hat\beta_{\rm full})\big]
\;\ge\; \frac{4\sigma^4}{k^2}\,\lambda_F^{(R)} \;>\; 0,
\end{equation}
and the gap decomposes additively into a term growing with
$\lambda_F^{(R)}$, a second-order leverage term, the in-fit
contraction term, and the in-fit F-leverage of the full fit.
The identity is a forget-side statement. Test-side prediction
variance needs its own condition: under the whitened leverage
$\delta = \|G_R^{-1/2} X_F^\top X_F G_R^{-1/2}\| = O(k/n)$ the
relative change is bounded,
\begin{equation}
\label{eq:whiten}
0 \le v_R(x_\ast) - v_{\rm full}(x_\ast) \le
\delta\,v_{\rm full}(x_\ast),
\end{equation}
but
the unwhitened ratio $\|X_F^\top X_F\|/\|X_R^\top X_R\|$ does
not control it---for a fixed test point the two variances can
differ by a constant factor while $k/n \to 0$ (Supplementary
Appendix). Test-side flatness in our experiments is an
empirical observation, not a corollary of the identity.
\end{proposition}

\noindent The proposition supplies the forget-side signature
the measurement shows---forget-side statistics inflate under
retraining---and identifies the driver as the leverage mass
of the deleted block (plus the in-fit F-leverage of the full
fit), not the identity of the statistic. It is a mechanism
for the linear model; we do not claim it quantitatively
predicts the deep-network ratios measured in
Section~\ref{sec:eval}. The inflation is not an artefact of
ensemble size or of one class: Section~\ref{sec:eval}
measures it at every tested scale and across the three tested
classes, with flat test-side ratios (Supplementary Tables
S33, S34), while \S\ref{sec:allowance} shows it is not the
room the encoder spends.

\subsection{Overview}
ECLIPSE (\emph{Encoding that Crosses audIt Lines while Preserving
Signals of Erasure-eligible records}) is a provider-side
construction with two active modules and one readout protocol
(Fig.~\ref{fig:pipeline}). In the first stage, FIDE optimizes
under live acceptance, so every iterate is feasible; in the
second, after the accepted model is deployed, CAL-Z reads
individual deleted records. Section~\ref{sec:region} explains
a source of forget-side dispersion in the stated linear
model; the resource used by the tested encoder is identified
separately---boundary measurements and the dose response in
Section~\ref{sec:eval} show the role of utility-side
tolerance:

\textbf{FIDE (Feasible-Direction Encoding).} A contrastive
encoder that minimizes a forget-vs-construction objective
while keeping every committed iterate feasible
(Remark~\ref{rem:anytime}; Algorithm~\ref{alg:fdce}).

\textbf{Observed outcome: boundary pinning.} The guarantee is
feasibility only, $s(\theta') \le \hat q$; in the primary runs
every accepted final also sits at the boundary \emph{exactly},
through the utility metric $m_1$
(Section~\ref{sec:pinning}).

\textbf{CAL-Z (Readout Protocol).} After the model passes and
is deployed, a reference-equipped party reads individual
deleted records at working-point thresholds
(Section~\ref{sec:calz}). CAL-Z is an evaluation protocol, not
part of the provider-side construction.

\input{fig_pipeline}

\subsection{FIDE: Feasible-Direction Contrastive Encoding}
\label{sec:fide}

\paragraph*{Design rationale: the contrastive structure.}
The provider must make per-record membership of $F$ readable
downstream \emph{while} satisfying the acceptance constraint.
Three natural templates fail: plain forget-set fitting
delivers $+0.05$--$0.08$ against $+0.14$--$0.20$ for the
contrastive objective (the fitted records are not
\emph{separated} from same-class neighbors, so scoring cannot
rank them); retention training under the audit is
selection-dominated (Section~\ref{sec:wash}); and cross-class
encoding moves the non-monitored population without
per-record resolution. The contrastive structure outperforms
the tested one-sided alternatives, and its two terms serve
distinct roles: the $\mathbb{E}_F \ell$ term lowers the
forget-set loss; the $\lambda\,\mathbb{E}_{N_c}\ell$ term
raises the loss on same-class retained members, creating the
separation that downstream scoring exploits. Without $\lambda$,
the encoder fits $F$ but does not separate it from $N_c$;
without the $F$ term, the model diverges but carries no signal
about which records were deleted. In the reported sweep, $\lambda{=}0.5$ is the best tested
value: $\lambda{=}2$ eliminates the measured gains, and
$\lambda{=}0.5$ is the only swept value delivering CI-positive
gains on all three selected signal-bearing starts
(Supplementary Table S24); the sweep does not establish
universal necessity.

\paragraph*{Why feasible-direction descent.} The constraint
region is non-convex and non-smooth ($s$ is a maximum of six
standardized deviations), and two of its components live on
discrete grids. We do not claim these features make
projection or Lagrangian methods impossible; they make exact
acceptance hard to guarantee from those generic recipes. A
score-checked backtracking step---accept a candidate only if
the full audit score stays within threshold---retains a
passing iterate by construction, which is the property the
provider actually needs (Remark~\ref{rem:anytime}).

\paragraph*{Objective.} Minimize
\begin{equation}
J_{\rm CE}(\theta) = \mathbb{E}_{F}\ell(\theta;x,y) -
\lambda\,\mathbb{E}_{N_c}\ell(\theta;x,y)
\label{eq:ce-objective}
\end{equation}
by feasible-direction descent (Algorithm~\ref{alg:fdce}).

\paragraph*{The loop in words.} Each iteration draws a fresh
batch $B$ of 60 construction-pool members, computes the
objective gradient on $F \cup B$, proposes a candidate
$\theta_{\rm cand} = \theta - \eta\,\nabla J$ (step size
$\eta$), and submits it to
the acceptance test: the \emph{full} six-metric audit score is
evaluated in full (the six-metric score $s_6$ of
Eq.~(\ref{eq:a6}) over $m_1$--$m_6$), and the candidate is
accepted iff
$s(\theta_{\rm cand}) \le \hat q_{0.20}$ \emph{and} the
objective improves. On rejection the step is halved (up to 8
backtracks); on acceptance the step grows by $1.1\times$
(capped at $3\times10^{-3}$); 12 consecutive rejected
iterations stall-stop the loop. The audit evaluation dominates
the cost, and it is exact: no surrogate or per-metric shortcut
is used, which is what makes Remark~\ref{rem:anytime} hold.

\paragraph*{Gradient variant.} The double-backward objective
$J_{\rm G} = \mathbb{E}_F\|\nabla_x\ell\| -
0.5\,\mathbb{E}_{N_c}\|\nabla_x\ell\|$. On some starts the
gradient variant is stronger; on others CE wins---the capability
is the task-plus-constraint, and first-order methods carry it.

\paragraph*{Complexity.} Algorithm~\ref{alg:fdce} runs $T$
iterations (typically 60). Each iteration needs one
forward-backward pass on $|F| + |B|$ records plus one full
audit evaluation (all six metrics; backtracking can invoke
this scorer repeatedly). The audit evaluation dominates the
cost; end-to-end, an encoding takes approximately $2$--$4$
minutes on a single consumer GPU.

\input{alg_fdce}

\subsection{Boundary Pinning as a Designed Outcome}
\label{sec:pinning}

\begin{definition}[Pinned final]
\label{def:pinned}
A returned model $\theta'$ is \emph{pinned} iff
$s(\theta') = \hat q$ exactly and the binding component of the
joint score is the utility metric: $\arg\max_i
|m_i(\theta')-c_i|/\sigma_i = 1$ (the $m_1$ term).
\end{definition}

Pinning is an observed endpoint property of the primary
accepted finals: all eleven sit at $s_6 = \hat q$ through
$m_1$, recording the same boundary value $m_1 = 0.7814$
(Section~\ref{sec:eval}), and the property recurs across
starts, requests, and audit instances. A plausible mechanism,
offered as a hypothesis rather than a law: the accuracy-type
metrics ($m_1$ overall, $m_2$ class-level) live on discrete
grids whose step is one evaluation-set flip, and the observed
trajectories exhaust utility-grid slack before the forget-side
metrics leave their bands, so descent ends pinned through
$m_1$. We have not proven monotone utility degradation along
accepted steps, and Definition~\ref{def:pinned} is not claimed
to hold for every construction or audit variant (the
expansion includes non-pinned, non-passing endpoints).

Two consequences follow. First, the audit budget is spent on
utility, not on the forget-side instruments: the audit's own
membership statistic stays within band on every measured final
($z_5 \le 2.19$ against the threshold $5.245$,
Table~\ref{tab:mb-paired})---this is the property that keeps
the encoding invisible to the suite's privacy instrument.
Second, pinned finals are brittle to threshold tightening but
not to honest coverage: a $1\%$ tightening
($0.99\hat q$) rejects all pinned finals while leaving held-out
honest acceptance unchanged (Section~\ref{sec:eval}); an
encoder re-run under the tightened threshold lands at the
nearest feasible grid value just below the new boundary
(Section~\ref{sec:eval}), not necessarily at an exact pin.

\subsection{CAL-Z: Reference-Calibrated Readout Protocol}
\label{sec:calz}
After the encoded model passes the audit and is deployed, a
reference-equipped party reads individual deleted records.
The protocol has four frozen steps.
(1)~\emph{Score}: each candidate record $i$ receives a score
$r_i$ under the deployed model (any scorer whose direction is
fixed in advance; the evaluations of this paper use
$r_\theta = -\ell$ and $-\|\nabla_x \ell\|$).
(2)~\emph{Per-record calibration}:
\begin{equation}
\label{eq:calz}
z_i = \frac{r_i - \mu_i}{\sigma^z_i},
\end{equation}
where $\mu_i$, $\sigma^z_i$ are the per-record mean
and standard deviation across the audit's own honest reference
ensemble---per-record difficulty calibration
following~\cite{carlini2022membership}; easy and hard records
are standardized against their own reference behavior, which is
what separates ``unusually easy for this record'' from
``member-like''.
(3)~\emph{Threshold freezing}: the working-point threshold is
fixed on a never-trained auxiliary set $N_{\rm aux}$ at a
nominal FPR, before any target record is scored.
(4)~\emph{Evaluation}: decisions on the target records are
reported with the actual TPR and actual FPR measured on the
disjoint never-trained set $N_{\rm eval}$.
Our use of per-record calibration as a readout of models
inside a live audit's acceptance region is the application,
not a contribution. The protocol requires reference-ensemble
access (Tier-2 capability, Section~\ref{sec:threat});
CAL-Z is an evaluation protocol, not part of the provider-side
construction.

\subsection{Design Choices and Their Validation}
\label{sec:design-choices}
Each frozen design choice of ECLIPSE is either validated by a
preregistered ablation or is structural to the problem;
Table~\ref{tab:design} maps them. Two choices deliberately do
\emph{not} carry weight: the iteration budget (gains saturate
at 15 iterations) and the deletion-set selection rule (a
uniformly random set with a fresh audit instance replicates),
which is why neither is defended.

\begin{table}[t]
\caption{Frozen design choices of ECLIPSE and where each is
validated. ``Load-bearing'' = removing or changing the choice
eliminates the capability in the cited measurement.}
\label{tab:design}
\centering
\footnotesize
\setlength{\tabcolsep}{3pt}
\begin{tabular}{p{24mm}p{14mm}p{30mm}}
\toprule
Choice & Setting & Validated by \\
\midrule
Contrastive structure & two-term Eq.~(\ref{eq:ce-objective}) & plain/one-sided $+0.05$--$0.08$ vs.\ $+0.14$--$+0.20$ (\S\ref{sec:ablations}) \\
Contrast weight $\lambda$ & $0.5$ & sweep optimum; $\lambda{=}2$ kills; only $\lambda{=}0.5$ CI-positive on $3/3$ signal starts (S24) \\ 
In-loop exact feasibility & full $s(\theta)$ per step & anytime termination (Rem.~\ref{rem:anytime}); pinned finals (Def.~\ref{def:pinned}) \\
Construction pool $N_c$ & 300 records & gradient variant collapses at $N_c{=}30$ (independent eval., S24) \\
Evaluation resolution & 200 non-members & frozen size; readings replicate at this resolution (S24) \\
Readout calibration & per-record $z$ & working-point readout vs.\ raw scores (Table~\ref{tab:zreadout}) \\
\midrule
\multicolumn{3}{l}{Not defended: iteration budget (saturates at 15);} \\
\multicolumn{3}{l}{deletion-set selection rule (random set replicates, 4/6 CI-positive).} \\
\bottomrule
\end{tabular}
\end{table}

\section{Evaluation}\label{sec:eval}
This section evaluates the joint event that a constructed
model passes its audit and increases membership readability.
In the primary scoring execution, eleven of the twelve CE/GR
finals pass, and all eleven accepted finals sit exactly at the
acceptance threshold, pinned through the utility metric; the
expansion probes then separate acceptance from signal delivery
across reference families, classes, scales, and calibration
assumptions, and include non-passing finals (the primary GR
e4-s2 endpoint and one dog-class start). The evaluation
proceeds in three acts that answer the three research
questions: \emph{Act 1} establishes feasibility and the
delivered signal, probing robustness by ablations, a
design-space sweep, and an alternative wash construction;
\emph{Act 2} explains why the audit tolerates the encoding,
from the decode gap through the dispersion measurement to the
two-part mechanism---admissibility and the utility
allowance---that gates delivery; \emph{Act 3} evaluates the
auditor's countermeasures instrument by instrument.

\subsection{Experimental Setup}
\textbf{Datasets and models.} All experiments use CIFAR-10
($n{=}50{,}000$) with the standard
split; the primary architecture is ResNet-18, replicated on
MobileNetV2 and on CIFAR-100. The full shared configuration
is fixed in
Supplementary Table S43; per-study deviations are listed
in the experiment manifests (Supplementary Table S19) and the
archived preregistration cards. \textbf{Reference families.} A preregistered probe with
exact-retrain references (15 fresh from-scratch retrains)
shows acceptance survives there (6/6 finals pinned and
passing) while the per-record signal does not for either
construction---first-order gains near zero, second-order
gains negative (Supplementary Table S23)---so the
fine-tuning-referenced scope is a measured boundary observed
for both tested constructions. \textbf{Wash statistics.} Wash-specific
statistics
($L$ selection ceiling, $W$ wash elevation, $G = W - L$,
$\gamma$) are defined with the supplementary wash tables;
$I_1$ is Definition~\ref{def:i1}.

\textbf{Determinism.} Base models $\theta_0$ and honest retrains
$\theta_R^{(s)}$ use the identical recipe of
Supplementary Table S43 under full deterministic seeding
(the retrains on $D \setminus F$). Audit metrics are computed
in the records' canonical (unaugmented) form, matching the
training input representation. Unless stated, ensembles have
three retraining seeds; the retrial study uses six.

\textbf{Protocol pools.} All data draws are within the deleted class, frozen before
each run: construction pool $N_c$ (300 records),
seventh-statistic set $N_a$ (100), threshold-calibration set
$N_{\rm aux}$ (100), and evaluation set $N_{\rm eval}$ (200)
are mutually disjoint. The primary $A_6$ split is six fit,
twelve calibration, and twelve held-out honest models at
$\alpha{=}0.20$; all three are one-epoch fine-tunes from
$\theta_0$ on the retained set. The encoding optimizer's
frozen hyperparameters are in Supplementary Table S43
($LR_0{=}10^{-3}$ in Algorithm~\ref{alg:fdce}).
Both materiality tiers ($I_1 \ge 0.10$ / $\ge 0.20$) are
preregistered with three-state verdict rules.


\textbf{Uncertainty.} We report target-record resampling
sensitivity intervals (2{,}000 resamples, same indices for base
and final, each model paired with its own non-member scores).
We do not claim coverage of non-member sampling, retraining, or
cross-request variability. Results are from a single preregistered
execution per configuration unless otherwise stated. The
intervals are within-execution statements, conditional on the
trained models, audit instance, and non-member sets; they do
not estimate end-to-end execution or retraining variance.
The observed cross-execution scoring differences (base
evaluations $\le 0.011$, Supplementary Table S6 note;
encoded-model downstream evaluations up to $0.014$,
Supplementary Table S18) are differences, not a variance
bound, and the intervals do not cover them; the record also
includes a larger, unreconciled discrepancy (an e4-s1 base AUC
of $0.545$ in the strong-probe run against $0.498$--$0.505$
elsewhere, Supplementary Table S27), which we report rather
than bound. Two of the nine encoding gains (e3-s4 $+0.027$,
n-s2 $+0.028$) are only about twice that noise, so their
individual positivity carries execution-level uncertainty;
the seven gate-clearing gains ($+0.102$ to $+0.246$) all sit
far beyond it.
\subsection{Act 1: Feasibility and Delivered Signal}
This act answers RQ1: does the encoding pass while carrying
readable per-record signal, and how far does that claim reach?

\subsubsection{Effectiveness of ECLIPSE: Label-Blind Results}
Scoring direction fixed in advance
($r_\theta = -\|\nabla_x\ell\|$, equivalently low-loss;
\S\ref{sec:decode-gap} shows output channels are
close on these finals). Thresholds for working points are frozen
on the auxiliary set ($N_{\rm aux}$, 100 records); we report
the actual TPR and actual FPR measured on $N_{\rm eval}$;
uncertainty statements follow the Setup convention.

\begin{table}[t]
\caption{Encoding results: label-blind AUC (base $\to$
encoded), target-record resampling sensitivity intervals of
the gain (member-index resampling only; training-seed and
non-member sampling variability are not covered). All finals
pass the frozen audit. Two deletion requests, two pools,
nine starts, first-order method.}
\label{tab:encoding-main}
\centering
\small
\setlength{\tabcolsep}{3pt}
\begin{tabular}{lccc}
\toprule
Start & base & encoded & interval \\
\midrule
$F$: e3-s1 & 0.547 & 0.703 & $[0.072, 0.251]$ \\
$F$: e3-s3 & 0.535 & 0.781 & $[0.160, 0.333]$ \\
$F$: e3-s4 & 0.517 & 0.544 & $[0.003, 0.055]$ \\
$F$: e4-s0 & 0.493 & 0.632 & $[0.073, 0.217]$ \\
$F$: e4-s1 & 0.505 & 0.607 & $[0.010, 0.189]$ \\
$F$: e4-s2 & 0.505 & 0.664 & $[0.087, 0.236]$ \\
$F'$: n-s1 & 0.596 & 0.737 & $[0.072, 0.217]$ \\
$F'$: n-s2 & 0.582 & 0.610 & $[0.001, 0.066]$ \\
$F'$: n-s3 & 0.650 & 0.790 & $[0.065, 0.215]$ \\
\midrule
\multicolumn{3}{l}{9/9 intervals exclude zero; 7/9 clear $0.05$
gate} & \\
\bottomrule
\end{tabular}
\end{table}

Table~\ref{tab:encoding-main}: nine of nine runs across two
deletion requests show independent label-blind gains with
per-run sensitivity intervals excluding zero; seven of nine
clear a $0.05$ working gate (gate on the point-estimate gain;
under the stricter interval-lower-bound reading, six of
nine)---we report both counts. The nested structure (six
starts under $F$, three under $F'$) is not a
$2{\times}2{\times}9$ cross; the new request provides one
preregistered task-level replication. Start selection: the
six $F$-pool starts were fixed in the preregistered protocol
before label-blind evaluation (weak-start readings retained,
not replaced; the designation record is with the experiment
manifests, Supplementary Table S19). All nine starts share
the original $\theta_0$ training run and the frozen record
sets; the two requests use their own audit instances (the
$F'$ instance has $\hat q{=}2.831$), so the starts span two
request-specific audits while sharing the remaining
infrastructure. The complete per-start listing is in
Supplementary Table S3.

\subsubsection{Impact: The Acceptance--Readout Gap}
\label{sec:zreadout}
The audit--decoder coexistence (Section~\ref{sec:decode-gap})
uses the audit's own $m_5$; here we report the readout's
operating points (protocol in Section~\ref{sec:calz}).

\begin{table}[t]
\caption{The acceptance--readout gap, per final:
reference-calibrated $z$-neggrad TPR at the nominal $5\%$
threshold (thresholds frozen on the auxiliary set), actual
FPR on the evaluation set (paired base$\to$final for both TPR
and FPR), and audit status ($\hat
q{=}5.245$). All eleven passing finals sit at the acceptance
value exactly, through the utility metric ($m_1{=}0.7814$);
GR e4-s2 sits one accuracy-grid step outside. Readings pair
with the primary scoring execution; three grid-boundary
finals are device-dependent in $A_6$ status (Supplementary
Table S12).}
\label{tab:zreadout}
\centering
\small
\setlength{\tabcolsep}{3pt}
\begin{tabular}{lcccc}
\toprule
Final & TPR b$\to$f & FPR b$\to$f & audit \\
\midrule
CE e3-s1 & 0.10$\to$0.33 & .115$\to$0.115 & $=\hat q$ \\
GR e3-s1 & 0.10$\to$0.43 & .115$\to$0.130 & $=\hat q$ \\
CE e3-s3 & 0.03$\to$0.23 & .100$\to$0.085 & $=\hat q$ \\
GR e3-s3 & 0.03$\to$0.10 & .100$\to$0.095 & $=\hat q$ \\
CE e3-s4 & 0.13$\to$0.07 & .070$\to$0.045 & $=\hat q$ \\
GR e3-s4 & 0.13$\to$0.07 & .070$\to$0.075 & $=\hat q$ \\
CE e4-s0 & 0.03$\to$0.23 & .065$\to$0.080 & $=\hat q$ \\
GR e4-s0 & 0.03$\to$0.17 & .065$\to$0.100 & $=\hat q$ \\
CE e4-s1 & 0.07$\to$0.13 & .035$\to$0.040 & $=\hat q$ \\
GR e4-s1 & 0.07$\to$0.17 & .035$\to$0.070 & $=\hat q$ \\
CE e4-s2 & 0.03$\to$0.23 & .070$\to$0.095 & $=\hat q$ \\
GR e4-s2 & 0.03$\to$\textbf{0.60} & .070$\to$0.070 & 5.258 \\
\bottomrule
\end{tabular}
\end{table}

Table~\ref{tab:zreadout} is the paper's central exhibit.
Nine of the eleven passing finals improve from honest-baseline
levels ($0.03$--$0.13$) to $0.10$--$0.43$ at actual FPRs of
$0.040$--$0.130$; the strongest (GR e3-s1) reaches TPR
$0.43$ at an actual FPR of $13.0\%$, where TPR resolution is
one record in $30$ ($0.43{=}13/30$, exact 95\% interval
$[0.26, 0.63]$; the boundary final's $0.60{=}18/30$,
$[0.41, 0.77]$). The median passing final reads $0.17$---the
maximum is a best case over twelve inspected endpoints, not a
typical attack. All eleven sit at the acceptance value exactly
($s_6{=}\hat q$), pinned through $m_1$ (all record
$m_1{=}0.7814$; Section~\ref{sec:pinning}); the twelfth
(GR e4-s2, $m_1{=}0.7813$, $s_6{=}5.258$) sits one grid step
outside and carries the strongest reading of all: TPR
$0.60$, and $0.47$ at the nominal $1\%$ threshold (actual
FPR $3.0\%$). Two passing finals (e3-s4, both families) show
no improvement; at the start level, five of six starts show
at least one improving passing final. What the table proves
is the gap itself: a model that satisfies every audited
aggregate to the letter can still carry per-record membership
signals readable at working-point thresholds. The full
per-final detail, including audit scores and a
boundary-mechanism recomputation, is in Supplementary Table
S7.
The boundary-pinned mechanism has a direct robustness reading
(zero-training recomputation on the cached scores): because
every passing final sits at exactly $s_6{=}\hat q$, a $1\%$
threshold tightening ($0.99\hat q$) rejects all eleven
passing finals while leaving held-out honest acceptance
unchanged ($8/12$; at $0.95\hat q$ honest acceptance drops to
$7/12$). But tightening is not a stable defense: re-running
the encoder \emph{under} the tightened threshold
(preregistered, three starts per level) converges just below
it ($s_6{=}5.182$ at $0.99\hat q$, $4.981$ at $0.95\hat q$)
and the signal survives on two of three starts per level
(point gains $+0.14$ to $+0.16$, CI lower bounds $+0.021$ to
$+0.038$; Supplementary Table S22). Threshold tightening
moves the race; it does not end it.

\emph{Architecture replication.} Under a freshly calibrated
$A_6$ instance on MobileNetV2 (30 archived honest references,
preregistered gates), all six encoded finals pass, four of
six label-blind AUC intervals exclude zero (up to
$0.509\to0.627$), and $z$-TPR@5\% improves on five of six
finals (maximum $0.23$ at an actual FPR of $7.5\%$; ResNet-18
tops at $0.43$; Supplementary Table S8)---the gap, and the
audit's blindness to it, are architecture-stable.

\emph{Dataset boundary, measured.} On CIFAR-100 (class pool
of 500 with adapted same-class sets drawn train-side), all
six encoded finals pass the fresh audit
instance---acceptance is dataset-stable---but the calibrated
readout gain does not replicate ($z$-TPR@5\% improves at most
$+0.07$ against a $0.20$ gate, and only three of six raw-AUC
intervals exclude zero; Supplementary Table S10); the raw
channel is additionally confounded by deletion-set selection
(honest bases already separate the selected records at AUC
$0.895$--$0.995$). A pool-geometry discriminator supports the
attribution: CIFAR-10 with the pool cut to 500 reproduces the
same dissociation (raw gains up to $0.56\to0.87$, calibrated
gain zero on every final; Supplementary Table S14); the
adaptation bundles several changes.

\emph{Class dependence, measured.} Across three further
deletion classes (frog, dog, airplane; 24 fresh honest
fine-tunes and a fresh $A_6$ instance each), acceptance is
class-stable while the per-record signal is not: frog
replicates strongly (4/6 starts CI-positive, gains $+0.10$
to $+0.24$), dog is partial (1/6), and airplane shows the
acceptance--signal separation (6/6 finals pass, no
CI-positive gain)---the same signature as the CIFAR-100
boundary, confirmed on independent audit instances (frog
6/6 CI-positive, airplane 0/6 again). The frog finals carry
ranking-level but not working-point signal ($z$-TPR
increments at most $+0.07$, washed after one epoch of benign
fine-tuning), so the Tier-2 readout is cat-specific among
the classes tested; a MobileNetV2 $\times$ dog crossing
attenuates but does not erase the signal. Within the
fine-tuning-reference family the class-specific thresholds
differ (frog $\hat q{=}16.8$ vs dog/airplane $4.5$);
separately calibrated quantiles do not by themselves
establish which class has wider physical tolerance, so we
record the transfer result without that inference. The
positive domain of the per-record claim is measured on 2/4
classes rather than assumed (Supplementary Table S23).

\emph{Scale boundary.} The encoding is demonstrated at
$k{=}30$ ($0.06\%$ of training data). At $k{=}100$ the
preregistered label-blind experiment (six starts) produced
no CI-positive gain ($+0.01$ to $+0.03$ straddling zero). At
$k{=}300$ an exploratory run found no calibrated-readout
replication ($z$-readout AUC $0.44$--$0.51$; raw ranking
unmeasured), and under an independently recalibrated
$k{=}300$ audit ($\hat q{=}5.503$ from 18 fresh honest
references) three encodings again pin and pass with no
signal delivery ($z$-TPR median $0.02$ vs honest $0.03$;
Supplementary Tables S16, S23, S35)---the boundary is not an
artifact of threshold reuse.

A controlled contrast bounds what the readout detects: three
natural non-compliance constructions (undertrained
unlearning, incorrect learning rate, partial deletion) show
member-level increments at or below $+1.2$ pp (against a
$+10$ pp gate), and the incorrect-lr variant's 99.9\%
model-level detection rate is matched by a 100\% rate on
decoy-only control pools---the readout detects a
\emph{model-level} signature for this family, not a
record-level membership signal. The encoding's gradient
updates also shift non-monitored classes (dog-class median
$z$-shift $-1.85\sigma$, ship $-0.79\sigma$)---a detection
opportunity for the auditor, not a per-record signal.


\subsubsection{Scope and Boundaries of the Demonstrated Harm}
\label{sec:working-points}
Two readouts, two conclusions. Under the fixed decoder
(thresholds frozen on the auxiliary set at nominal 5\% FPR;
actual FPR $0.020$--$0.060$), TPR@5\% is $\le 0.10$: a
label-blind outside party realizes \emph{ranking-level}
membership inference, not confident individual attribution
(a $0.10$ TPR at $0.05$ nominal FPR with a $1\%$ membership
prior gives a ${\sim}2\%$ positive predictive value). The
reference-calibrated readout is stronger (TPR up to $0.43$
among passing finals, $0.60$ boundary-adjacent; PPV about
$3\%$ at the same prior) but conditional on
reference-ensemble access. Low-FPR identification of
individual records is not established by either reading;
what is established is the gap itself---models that pass
every audited statistic carry readable per-record signals
(consolidated harm matrix: Supplementary Table S17). A
persistence probe closes the scope: after one epoch of
benign fine-tuning on retained data, the median paired
retention of the $z$-TPR gain is $0.0$ (range $-2.5$ to
$1.0$; three epochs, $0.14$; Supplementary Table S9)---the
demonstrated readability is a property of the returned
artifact at the audit interface, not a persistent watermark
under continued training.

\emph{Sequential requests, probed on two pairs.} A second
deletion request $F_2$ (disjoint, same class) arriving at a
passing encoded final was probed end to end on two
independent request pairs. All three signatures replicate
on both: every stale model is accepted under the fresh
request-2 audit (pair 1's stalest scores $6.25$ against
$\hat q{=}6.53$)---sequential protection inherits the
reference-operator's tolerance; honest processing of $F_2$
reduces the first request's signal ($z$-TPR $0.17\to
0.00$--$0.03$ on pair 1, $0.47\to 0.07$--$0.23$ on pair
2); and the gap does not compound---on pair 1 the second
encoding re-delivers on $F_2$ ($0.00\to0.23$) with $F_1$
near baseline, while on pair 2 the passing second encoding
delivers no working-point gain on $F_2$
(start-dependence, as in the primary experiments). The gap
is per-request, not cumulative (Supplementary Tables S11,
S31).

Deleted-vs-retained AUC (raw negative-loss, $F$ vs the
3000-record retained sub-sample) improves on the strongest
gradient-variant starts ($0.27\to0.60$ on e3-s3,
$0.17\to0.60$ on e4-s0; per-start table in Supplementary
Table S21). $F$-vs-retained and $F$-vs-never-trained are
different inference tasks; we report both without aggregating
across them.


\subsubsection{Ablations}\label{sec:ablations}
\textbf{Protocol requirements, established empirically.}
Construction-set diversity ($N_c{=}300$, required by the
gradient variant; CE replicates at 30) and evaluation
resolution ($200$ non-members) are protocol requirements,
not tunables: at $N_c{=}30$ the gradient variant collapses
under independent evaluation ($0.49$--$0.59$;
$30$-vs-$30$ swings ${\pm}2\sigma$). No universal minimum is
established; the small-pool failure belongs to the gradient
protocol we tested. A preregistered ablation of the frozen
encoder choices (Supplementary Table S24) locates which of
them are load-bearing: the contrastive weight
$\lambda{=}0.5$ is critical and sits at the sweep optimum
($\lambda{=}2$ kills the attack; only $\lambda{=}0.5$
delivers CI-positive gains on all three signal starts), the
construction-pool size and iteration budget are not (CE
replicates down to $N_c{=}30$; 15 iterations saturate), and
the deletion-set selection rule does not gate the attack (a
uniformly random 30-record deletion set with its own fresh
audit instance replicates at 4/6 CI-positive).


\paragraph*{Construction ablation}
Same starts, budget, and audit; all arms evaluated by the same
label-blind protocol (construction-set evaluation; the
independent-evaluation gains of
Table~\ref{tab:encoding-main} span $+0.03$ to $+0.25$).
Contrastive first-order CE (Eq.~%
\ref{eq:ce-objective}) reaches $+0.14$--$+0.20$ (best start);
plain $F$-CE is weaker; one-sided gradient objectives
($F$-only $+0.05$--$0.08$; $N_c$-only ${\approx}0$) show the
joint separation is what transfers; class-bias grids are
near-null. The contrastive structure, not a particular loss
form, is what carries the capability.


\subsubsection{Design-Space Observation under Intersection Constraints}
A five-alternative exploration was run under a stricter
acceptance object than either the marginal or the
single-instance joint rule: a \emph{multi-instance
intersection constraint} from $K{=}10$ bootstrap audit
instances, where a candidate is accepted only if its
conformity ratio is within threshold on all $K$ instances
simultaneously. Under this constraint, among six
systematically tested constructions only the double-backward
gradient objective produced CI-positive label-blind gains
(design starts $4/6$, prospective $3/6$, MobileNetV2 $2/6$);
the five first-order alternatives (CE-loss regularization,
leading-eigenvector control, adaptive signal--robustness
weighting, layer-selective gradients, cross-class encoding
via non-monitored logits) produced accepted finals on most
starts ($4$--$6$ of $6$) but no CI-positive start. This is an
observation under one constraint family, not a necessity
claim, and two of our own measurements bound it: under the
single-instance joint rule $A_8$, five of six
marginally-constructed first-order CE finals are accepted
(bootstrap-fragile; Supplementary Table S12), and a
first-order CE encoder optimizing directly under $A_8$ passes
on every tested start ($12/12$)---non-inferior where it
fires, start-dependent where it does not (preregistered;
Supplementary Table S12). We offer no geometric law behind
the intersection result.

The signal capability is start-dependent (${\sim}50\%$ of
honest starts produce signal-carrying finals), and no
geometric criterion we tested predicts which starts succeed;
under a \emph{conditional i.i.d.} projection a provider
attempting $N$ starts reaches $1-(1-0.5)^N$ ($>96\%$ at
$N{=}6$), with wide inherited uncertainty (exact 95\%
interval $[0.26, 0.74]$) and possible cross-start
correlation; the pooled per-start rate is $9/18$ CI-positive
(Supplementary Table S15), while the canonical CE attack's
own starts (Table~\ref{tab:encoding-main}) are $9/9$. The
encoding trajectory is non-monotonic, so start selection
requires full encoding runs rather than cheap early probes.



\needspace{3\baselineskip}
\subsubsection{The Mean-Objective Wash: Region Geometry and
Recovery-Score Controls}\label{sec:wash}
\begin{definition}[Normalized forget-margin recovery]\label{def:i1}
For a returned model $\theta'$ with paired honest reference
$\theta_R$ (the same unit's retrain that $\theta'$ was derived from),
\begin{equation}
\label{eq:i1}
I_1(\theta') = \frac{\mathrm{fm}(\theta') - \mathrm{fm}(\theta_R)}
{\mathrm{fm}(\theta_0) - \mathrm{fm}(\theta_R)},
\end{equation}
where
$\mathrm{fm}$ is the mean correct-class logit margin on $F$. The
denominator (the paired deletion delta) is positive in every unit we
report. $I_1$ is a normalized behavioral differential---a recovery of
the forget-margin delta---not a measure of information content or
privacy harm; in particular it is also computable for honest models
via cross-seed comparison, which is exactly the honest-seed-luck
control of the controls paragraph below. A wash is feasible iff
$\theta' \in \mathcal{A}$; the preregistered engineering tiers
(\emph{material} at $I_1 \ge 0.10$, \emph{strong} at $I_1 \ge 0.20$)
refer to this proxy only.
\end{definition}
The region the encoding exploits was mapped with a simpler
object, the \emph{mean-objective wash}, which maximizes the
normalized forget-margin recovery $I_1$ under the banded audit. Rollback is
excluded on every tested path ($0/240$ interpolation points
across twelve units and three scales; greedy per-layer
rollback best $I_1 = -0.03$; Supplementary Table S20), so the
audit is not vacuous. The interior admits unit maxima
$I_1^* = 0.179$--$0.326$ at $k{=}30$ and $0.142$--$0.214$ at
$k{=}300$, first-step dominated (Supplementary Tables S1--S2),
replicates on CIFAR-100, and is unchanged by stronger audit
statistics (RBF-MMD, an MLP-based MIA, per-sample
extreme-value tests) or by removing any single metric---the
binding constraint behaves like the region's width, not the
particular statistics (Supplementary Table S21). Banding the
forget-margin directly shows the elevation cap depends on
which band shape the auditor chooses: a one-sided empirical
cap rejects five of six fine-tuned washes and $37/62$
cheap-manipulation points where a symmetric band barely
touches them (Supplementary Table S40). Finally, the controls
reinterpret recovery scores: at exact-retrain references,
matched-budget selection accounts for all but
$0.000$--$0.078$ of the recovered margin across five six-seed
units (all below the $0.10$ gate), a one-scalar logit-scale
control is equally deflating, and the single gate-clearing
ft-referenced gain ($+0.158$) fails the frozen
cross-ensemble rule---margin-recovery readings without these
controls answer questions that honest processing already
resolves. Approximate-operator references widen the region
for the same reason ($I_1^*$ to $0.50$/$0.40$ at the two
scales).
\label{sec:fmband}\label{sec:selection-control}


\subsection{Act 2: Why the Audit Tolerates the Encoding}
This act answers RQ2: it locates the room the audit grants,
measures where the encoder's budget goes, and separates
admissibility from the allowance.

\subsubsection{Why the Audit Does Not Detect the Encoding}
\label{sec:decode-gap}
The dispersion argument of \S\ref{sec:region}
predicts the audit has room; this subsection measures where
the encoder's budget actually goes. On the six
gradient-variant finals, member-vs-fresh AUC under
$-\ell$, margin, confidence, and $-\|\nabla_x\ell\|$ are close:
maximum inter-channel difference 0.016 (strongest final:
$0.810/0.813/0.810/0.818$; per-final detail in Supplementary
Table S5). This proximity does not
prove information equivalence across scorers, nor does it
automatically describe the nine CE finals. The encoding is
an \emph{output-space} phenomenon on these finals.

\begin{table}[t]
\caption{Audit--decoder gap: per-final paired values. Base =
audit-passing honest start; Final = encoded model (gradient
variant). $A_{\rm down}$ = downstream $-\|\nabla_x\ell\|$ AUC vs.\ 200
fresh non-members; $m_5$ = audit's membership statistic;
$z_5$ = $|m_5 - c_5|/\sigma_5$ (final). Five finals sit at
the $A_6$ acceptance value ($s_6{=}\hat q{=}5.245$, pinned
through the utility component); e4-s2 measures
$s_6{=}5.258$, marginally above (\S\ref{sec:zreadout}).}
\label{tab:mb-paired}
\centering
\small
\setlength{\tabcolsep}{3pt}
\begin{tabular}{lcccccc}
\toprule
 & \multicolumn{2}{c}{$A_{\rm down}$} & \multicolumn{2}{c}{$m_5$} &
$z_5$ & $s_6$ \\
Start & base & fin & base & fin & fin & fin \\
\midrule
e3-s1 & .547 & .714 & .661 & .737 & 0.19 & $=\hat q$ \\
e3-s3 & .535 & .562 & .719 & .607 & 2.14 & $=\hat q$ \\
e3-s4 & .517 & .543 & .734 & .604 & 2.19 & $=\hat q$ \\
e4-s0 & .493 & .631 & .820 & .611 & 2.08 & $=\hat q$ \\
e4-s1 & .505 & .626 & .784 & .755 & 0.08 & $=\hat q$ \\
e4-s2 & .505 & .818 & .814 & .731 & 0.28 & 5.258 \\
\bottomrule
\end{tabular}
\end{table}

Table~\ref{tab:mb-paired} shows the coexistence: downstream
ranking rises ($+0.03$ to $+0.31$) while the audit's $m_5$
deviation stays at most $2.19$ MAD against a $5.245$
threshold, with $m_5$ \emph{decreasing} on five of six finals (the exception is e3-s1, whose $m_5$ rises from $0.661$ to $0.737$). The joint score of e4-s2 crosses the threshold through the utility component $m_1$, not $m_5$ (Section~\ref{sec:zreadout}). Note: Table~\ref{tab:mb-paired} reports the gradient-variant finals; the nine CE finals of Table~\ref{tab:encoding-main} are a different set, and this analysis does not automatically describe them (the CE-final counterpart measurements are in Supplementary Table S18).
The gap is compatible with two dampers:
(i) a scorer/pool mismatch (logistic probe on two features
against a mixed-class retain pool vs.\ raw loss ranking
against same-class fresh non-members), and (ii) the tolerance
width itself. The study reports the coexistence on these six
paired observations. Both suggested dampers are now measured
rather than open: the pool half is separated (the calibrated
signal is stronger on the audit's own retained pool than on
fresh records, Supplementary Table S37), and halving the
forget-side tolerance width retains about $78\%$ of the
delivered gain (the forget arm at $\kappa{=}0.5$:
$0.121/0.155$; Supplementary Table S38), so the width
contributes but does not dominate. A general AUC-to-MAD
mapping remains open.
The per-metric acceptance rule, not overall displacement,
decides acceptance.


\subsubsection{The dispersion measurement}
\label{sec:dispmeas}
Quantitatively (Table~\ref{tab:calibration}), the honest ensemble's
dispersion divided by the full-data band is $4.78/1.91/2.84$ for the
forget-loss metric at $k = 30/300/1000$ (small-sample statistics:
three honest seeds per scale; the non-monotone $1.91\to2.84$
movement is within what three-seed maxima produce, and the
qualitative pattern---forget-side inflation with test-side
flatness---is what the 40-retrain archival control supports;
see the mechanism paragraph below for the $n{=}6$
confirmation),
while the KS metric
lies above 1 at $k{=}30$ and marginally at $k{=}300$ ($1.05$), the MIA
metric only at $k{=}30$, and test-side accuracy dispersion
stays inside the band at all $k$ (a 40-retrain archival study
confirms test-side dispersion is flat in $k$). The pattern is
informative: aggregate fit statistics on the deleted records are
where constraint release bites hardest, separability statistics
less so, and untouched-population statistics not at all. The
calibration lesson is metric-selective and
scale-coupled: \emph{forget-fit statistics must be re-calibrated
at each deletion scale}. The ensemble band we use throughout is
the minimal correction and covers every honest model in the
calibration ensemble by construction; it is a sample statistic, not
a coverage guarantee---a deployed audit should fix a target
false-positive rate and verify coverage on fresh honest outputs,
and per-scale re-estimation is one sound way, not the only
defensible one, to get there.

The inflation is not an artefact of ensemble size or of one
class: at $n{=}6$ retrains the ratios become $7.5/2.1/3.9$ with
the signature unchanged (Supplementary Table S34), and the
class exhibit measures $1.5\times$, $2.7\times$, and
$18.2\times$ on frog, dog, and airplane with flat test-side
ratios (Supplementary Table S33). The phenomenon therefore
holds at three levels---exact linear-model mechanism,
scale-robust measurement, replication across the tested
classes---while, as \S\ref{sec:allowance} shows, it is not
what the encoder spends.


\subsubsection{What gates delivery: admissibility and budget}
\label{sec:allowance}

Two mechanisms, not one, decide whether an encoding can be
delivered: the region is \emph{centred on the reference
operator's forget-fit} (which decides whose history is
admissible), and within the admissible family the
\emph{utility allowance}---the slack the accuracy-side bands
($m_1$, $m_2$) leave below honest accuracy---sets how much
signal can be bought. A dose-response separates them.

\emph{Admissibility: under a retrain-calibrated audit, both
F-aware profiles are inadmissible.} Under that audit the bands are narrow
on every component (accuracy $\pm0.73$ points, forget-set loss
$\pm0.69$ around a centre of $0.82$, forget margin $\pm0.91$
around $3.28$; $\hat q{=}5.638$). Two candidate starts fall
outside, on different components. The original model $\theta_0$
is in band on accuracy ($z_1=-0.9$) and outside on the
statistics that record the deletion---forget-set loss $0.10$
against a centre of $0.82$ ($z_3=-5.8$) and forget margin
$5.33$ ($z_6=+12.8$, $2.3\times$ the band). The one-epoch
fine-tune start is outside on accuracy ($z_1=-55$) and on the
margin in the opposite direction ($z_6=-14$) while staying in
band on the loss ($z_3=+1.4$). Both candidate profiles are
therefore rejected, on different components, at comparable
distance on the margin (per-component z-profile in
Supplementary Table S41).

Reference family and start ancestry are collinear here---a
retrain reference implies a start that never saw $F$, a
fine-tune reference implies one that did---so the retrain-side
closure cannot be attributed to either alone. What can be said
is that it closes twice: neither tested start that saw $F$ is
admissible under such an audit, and the admitted retrain
starts deliver no positive signal under the tested
constructions (Section~\ref{sec:eval}). Their $0.73$-point
accuracy half-width lies below the dose-response study's
tested range, so that study does not independently explain
the retrain-side failure. We report the two closings jointly
rather than as a single identified cause---not an identified
operator-family effect.

\emph{Budget: what the encoder spends once admitted.} Inside
the fine-tune-referenced family the finals sit on the lower
edge of the accuracy component exactly: $m_1 = 0.7814$ against
$c_1 = 0.8231$ and $\sigma_1 = 0.00795$ gives
$z_1 = -5.245 = -\hat q$. The dose-response scales one
family's half-widths at a time: acceptance holds iff
\begin{equation}
\label{eq:dose}
|m_i(\theta) - c_i| \le
\kappa_{\mathrm{fam}(i)}\,\hat q_\alpha\,\sigma_i
\quad\text{for all } i,
\end{equation}
where $\kappa_{\mathrm{fam}(i)}{=}\kappa_U$ for the accuracy
components $i\in\{1,2\}$ and $\kappa_F$ otherwise, with
$(\kappa_U,\kappa_F)=(\kappa,1)$ in the utility arm and
$(1,\kappa)$ in the forget-side arm. This moves the mean
delivered gain monotonically over a
six-fold range (one start decreases slightly between
$\kappa{=}2$ and 3, $0.247\to0.236$):
$0.084 \to 0.125 \to 0.154 \to 0.188 \to 0.225 \to 0.253$ for
$\kappa \in \{0.5, 0.75, 1, 1.5, 2, 3\}$ (per-start values and
per-cell feasibility in Supplementary Table S38; at
$\kappa{=}0.5$ and $0.75$ the tightened audit no longer admits
two of the three starts, so those two points are single-start
probes). Scaling only the forget-side half-widths
($m_3$--$m_6$) leaves the mean at $0.154$--$0.155$ for
$\kappa \ge 0.75$, with every start within $0.002$ of its own
$\kappa{=}1$ value; at $\kappa{=}0.5$ the mean falls to $0.121$
as one start collapses (to $0.060$) and a second drifts by
$0.009$. Within an admissible family, then, the deletion-inflated forget-side
tolerance is not the resource the channel consumes---a
negative result the dispersion story had excluded.

The allowance orders delivery \emph{within} the working family
($4.2$ accuracy points) and is \emph{not sufficient}: the
k=100 instance ($4.45$ points), airplane ($7.5$) and dog
($8.1$) all lie above the working allowance and deliver nothing
or one start (Supplementary Tables S36, S39). Admissibility and the allowance together close the retrain
family; neither orders the class or the scale axis, which
therefore remain empirical boundaries rather than predicted
ones.

The readout is not an artifact of the comparison population
either: on the audit's own retained pool the calibrated signal
is stronger than on fresh records (median paired $z$-gain
$0.312$, cluster-bootstrap $95\%$ CI $[0.137, 0.377]$, against
$0.151$ $[0.033, 0.297]$; Supplementary Table S37), so the gap
is not eliminated by aligning the comparison population; how
much of it belongs to the instrument's shape, as opposed to
other readout-design choices, is not isolated here.

\begin{table}[t]
\caption{Calibration measurement (ResNet-18). Top: honest
unlearned-ensemble dispersion (three seeds per scale) divided
by the full-data retraining band (six retrains), per metric and
deletion size ($>$\,1 means the naive band under-covers the honest
ensemble). Middle: honest false-positive rate under the naive band.
Bottom: archival test-side control (40 full retrains at
$k \in \{0,500,2500,5000\}$; accuracy dispersion flat in $k$).}
\label{tab:calibration}
\centering
\small
\setlength{\tabcolsep}{4pt}
\begin{tabular}{lccc}
\toprule
Dispersion ratio vs.\ full-data band & $k{=}30$ & $k{=}300$ & $k{=}1000$ \\
\midrule
$m_3$ forget loss      & \textbf{4.78} & \textbf{1.91} & \textbf{2.84} \\
$m_4$ KS distance      & 1.58 & 1.05 & 0.79 \\
$m_5$ MIA AUC          & 1.02 & 0.56 & 0.29 \\
forget margin          & 1.24 & 0.40 & 0.32 \\
$m_1$ overall accuracy & 0.20 & 0.20 & 0.67 \\
\midrule
Honest rejections under naive band & $3/3$ & $2/3$ & $1/3$ \\
\midrule
\multicolumn{4}{l}{Test-side control: acc.\ std $0.0025 \to 0.0017$
(flat in $k$).}\\
\bottomrule
\end{tabular}
\end{table}

Inside the region, what binds differs by protocol: under
$A_5$ the forget-side metrics bind (MIA probe $m_5$ rides at
ratio $1.0$ to its band, forget-loss $0.92$--$0.95$, KS
$0.98$; the wash readings of Section~\ref{sec:wash}), while
under the $A_6$ encoding audit the binding component is
instead the utility metric (\S\ref{sec:zreadout}); neither
transfers as the other's mechanism.

\subsubsection{What the Audit Certifies: Calibration Evidence}\label{sec:approx-op}
Independent calibration (a fresh six-model max-deviation
band applied to a new ensemble) rejects $5/6$ honest
retrains: at deployable budgets the audit's binding
constraint is honest coverage, not attack resistance.

\emph{Fixed-request split (the guaranteed experiment).} On the
finetuning family (30 archived honest models on one request),
a fixed disjoint split designated before scoring (fit: 6,
calibrate: 12, test: 12) gives an honest rejection rate of
$4/12$ at $\alpha = 0.20$ (Clopper--Pearson 95\% CI
$0.10$--$0.65$) and $0/12$ at $\alpha = 0.10$. This is the
\emph{conditional} rate of one frozen audit; the conformal
guarantee is a \emph{marginal} statement over calibration
randomness, so the single result neither proves nor disproves
it. Resource accounting: max-deviation bands at the same
$6$-model budget reject $7/12$; matched-total-budget
max-deviation (18 references) rejects $2/12$. The honest
attribution: most of the coverage gain comes from using more
reference models; at matched budget the conformal rule is not
pointwise better---its contribution is the declared target
with a validity guarantee (re-split and rotation counts are in
the artifacts). \emph{Attacks under the same frozen rule}:
$5/21$ archived checkpoints rejected at $0.20$ (all three
calibration-aware washes among them; built against earlier
band audits, so not a full adaptive test), $0/21$ at
$0.10$. Honest and attack columns have different denominators
and generation processes; no detection comparison is implied.
The three rules side by side are in Supplementary Table S44.



\emph{Cross-request sensitivity (empirical stress test, no
guarantee claimed).} A pooled variant (centers from each
request's own retrains, scales pooled across the other
requests---adaptive and self-centered, hence outside the split
guarantee) applied to the original draw plus three
preregistered fresh draws rejects $1/12$, $2/3$, $0/3$, and
$0/3$ respectively (equal-weighted $0.19$; pooled $3/21$),
against $60/63$ for single-draw max-deviation bands under the
same shift: limited transfer in aggregate, heterogeneous
across draws. Accepting a model under this audit bounds false
rejection of honest outputs; it does not certify that any
accepted model has correctly unlearned.



Two formal objects anchor the empirical sections: the coverage
guarantee of the split-conformal audit, and the design-space
observation that bounds which encoding constructions carry signal
under multi-instance intersection constraints.

\subsubsection{Split-Conformal Calibration and Its Guarantee}
\label{sec:m3}
The coverage failure of honest ensembles
(Section~\ref{sec:approx-op}) is a calibration problem,
and calibration problems have standard tools. We calibrate the
banded audit with split conformal prediction
\cite{vovk2005algorithmic, angelopoulos2021gentle}.

\emph{The guaranteed procedure.} All
statements of guarantee in this paper attach to the fixed-request
split version and to nothing else:
disjoint honest sets $H_{\rm fit} \perp H_{\rm cal} \perp
H_{\rm test}$ on one request family; centers \emph{and} scales
frozen from $H_{\rm fit}$ alone; threshold $\hat q_\alpha$ from
$H_{\rm cal}$ alone ($\lceil (n_{\rm cal}{+}1)(1{-}\alpha) \rceil$-th
order statistic, infinite out of range, ties by order, $\sigma_i$
guarded at $10^{-12}$); accept iff
\begin{equation}
\label{eq:conf}
s(\theta) \le \hat q_\alpha.
\end{equation}
For a single scored output exchangeable with the calibration
scores given $H_{\rm fit}$, the marginal false-rejection
probability is at most $\alpha$---the standard split-conformal
statement, its probability over $(H_{\rm cal}, \theta_{\rm new})$
jointly. The unit of the guarantee is one scored output; rotated
splits and leave-one-ensemble rotations share calibration data and
are reported descriptively. The acceptance region of this procedure
is a \emph{third} object---its per-metric half-widths are
$\sigma_i \hat q_\alpha$, in general different from both the
symmetric max-deviation band and the one-sided empirical cap of
\S\ref{sec:fmband}; the three must not be merged into one defense.



\subsection{Act 3: Defense Evaluation}
\label{sec:mitigation}
This act answers RQ3. We evaluate four audit-side
instruments---the aligned seventh
statistic ($A_7$), the joint covariance rule ($A_8$), the
aggregate z-detector, and a stronger membership probe---with
threshold tightening as a fifth variant. The joint rule and
threshold tightening are measured against both frozen
constructions and preregistered adaptive re-optimization arms;
the aligned statistic is tested through re-optimization of the
unchanged objective and two objective redesigns; the
z-detector through a composite-constraint arm. The measured
benefit against a frozen construction need not persist under
re-optimization---except at $A_7$ and the z-detector, where
the compression survives every tested evasion tier.

\begin{table}[t]
\caption{A directly aligned statistic compresses the attack
under re-optimization. Paired runs: from the same three start
checkpoints and fixed record sets, the unchanged first-order CE
objective is optimized \emph{separately} under the $A_6$ and
the recalibrated $A_7$ audits; the table compares the two
resulting final models under the same evaluation and readout
(base AUC values from a separate base-execution; cross-run
difference $\le 0.011$, Supplementary Table S6 note). In this
subsection, \emph{re-optimization} (unchanged objective re-run
under a new rule) and \emph{objective redesign} are distinct;
this experiment is the former.}
\label{tab:mitigation-paired}
\centering
\small
\begin{tabular}{lcccc}
\toprule
 & base & gain$_{A_6}$ & gain$_{A_7}$ & pass$_{A_7}$ \\
\midrule
e4-s0 & 0.494 & +0.140 & +0.045 & \checkmark \\
e4-s1 & 0.498 & +0.150 & +0.038 & \checkmark \\
e4-s2 & 0.504 & +0.121 & +0.050 & \checkmark \\
\bottomrule
\end{tabular}
\end{table}

\paragraph*{The aligned seventh statistic ($A_7$).}
Table~\ref{tab:mitigation-paired} shows the paired compression:
adding $d(\theta)=\mathbb{E}_{N_a}\|\nabla_x\ell\| -
\mathbb{E}_F\|\nabla_x\ell\|$ (frozen on independent fit
models) to the joint score and recalibrating ($\hat q{=}6.30$)
drops the CE encoding gain from $+0.12$--$+0.15$ to
$+0.04$--$+0.05$ per start---the attack survives at roughly a
third of its strength. On the same batch of 12 honest models,
the seven-metric audit accepts $9$ where the base audit accepts
$8$ ($n{=}12$; one-model difference, no coverage claim;
artifact \texttt{gate-g3-20261007T141513}; paired gains in
Table~\ref{tab:mitigation-paired}, artifact
\texttt{mc-20261007T231754}). The
cost is compute: the gradient statistic requires per-record
backward passes. These gains are point estimates per start;
member-resampling intervals are not attached (the paired
per-start reductions of $0.07$--$0.11$ are far larger than the
cross-run base difference of at most $0.011$ disclosed in
Supplementary Table S6, but no interval column is computed).
This is a measured partial mitigation, not a
proven security--coverage frontier.

\paragraph*{Objective redesign at $A_7$.} Two redesigned objectives---an anti-statistic penalty that
neutralizes the seventh statistic directly, and the
second-order gradient-norm objective---were optimized under
the same recalibrated $A_7$ audit (same starts, record sets,
and evaluation; preregistered gates, Supplementary Table S26).
Neither recovers a CI-positive signal: the penalty arm is
compressed outright ($+0.007$ to $+0.016$), and the
second-order arm reaches $+0.040$ to $+0.062$ with every
interval crossing zero. The $A_7$ compression thus survives
both tested evasion tiers; broader redesign families remain
untested.

\paragraph*{A stronger membership probe.} Replacing the two-feature
logistic $m_5$ with an MLP probe on five output features (loss,
margin, entropy, true logit, top-other; all black-box-plus-losses
accessible) and recalibrating the audit on the same fit6/cal12/
test12 honest models yields a rule with identical honest
acceptance ($8/12$) that \emph{rejects all nine rescored primary
finals} (the six CE and three gradient-variant checkpoints;
their scores sit at the old threshold values, above
the recalibrated $5.220$---a threshold roughly $0.5\%$ below
the original, which alone rejects utility-pinned finals, so
the experiment establishes rejection by the upgraded rule
rather than detection gained from the MLP component. A fresh
encoding run
\emph{under} the stronger rule re-pins and delivers again:
CI-positive gains on three of three starts ($+0.162$, $+0.129$,
$+0.141$; all passing; Supplementary Table S27). The race
pattern of the rule-based defenses thus recurs when the
instrument and its calibration move together.

\paragraph*{The joint acceptance rule ($A_8$).}
\label{sec:a8}
The marginal suite scores each statistic against its own band.
A joint rule instead scores the six-metric vector against the
honest ensemble's covariance:
\begin{equation}
\label{eq:a8}
s_J(\theta) = \bigl(m(\theta)-c\bigr)^{\top}
\hat\Sigma^{-1}\bigl(m(\theta)-c\bigr),
\end{equation}
with $\hat\Sigma$ a shrunk
covariance of the fit ensemble ($\rho{=}0.5$, frozen before any
attack readout) and the threshold the conformal quantile of the
calibration ensemble's $s_J$ scores at the same $\alpha{=}0.20$.
The comparison is zero-training and uses the identical
fit/cal/test split as $A_6$.

\begin{table}[t]
\caption{Marginal ($A_6$) vs joint ($A_8$) acceptance at matched
$\alpha{=}0.20$ and identical reference split. The joint rule
accepts strictly more honest models and rejects the six
marginally-constructed gradient finals. The $A_6$ column is
from the primary scoring execution; three grid-boundary
finals are device-dependent (per-final record in
Supplementary Table S12).}
\label{tab:a8}
\centering
\small
\begin{tabular}{lccc}
\toprule
Rule & Honest (12) & Finals (12) & GR family (6) \\
\midrule
$A_6$ marginal max & $8$ & $11$ accepted & $5$ accepted \\
$A_8$ joint (cov.) & $12$ & $5$ accepted & $0$ accepted \\
\bottomrule
\end{tabular}
\end{table}

Table~\ref{tab:a8}: the joint rule accepts all twelve held-out
honest models---four of which the marginal rule rejects,
because their deviations are correlated, exactly the direction
a Mahalanobis form discounts---while rejecting all six
gradient-variant finals and one CE final at scores well beyond
the threshold ($s_J$ 17--25 against $\hat q_J{=}16.5$; five
CE finals pass at 12--14). The mechanism is the standard one
(honest co-movement lives in the top covariance directions,
the encoders' deviation pattern does not), and shrinkage
sensitivity confirms it directionally (at $\rho{=}0.75$
rejection falls to $5/12$, at $\rho{=}0.25$ it rises to
$10/12$, honest coverage $12/12$ throughout); the result is
one audit instance, twelve finals, and a frozen shrinkage
constant---a mechanism demonstration, not a deployed defense
(Supplementary Table S12). The closure is not
family-intrinsic: the anti-statistic and second-order
redesigns under $A_8$ add no evasion (no CI-positive gain on
any start, $+0.009$ to $+0.078$, all intervals crossing zero,
passing on all six runs; Supplementary Tables S12, S32),
whereas a first-order CE encoder passes $A_8$ on every tested
start ($12/12$ preregistered; CI-positive on $5/6$ design and
$2/6$ prospective starts) and the gradient-objective arm
passes with the strongest readout measured in this work (AUC
$0.85$ design, $0.82$ prospective; Supplementary Table S12).
The joint rule and threshold tightening each admit evasion by
re-optimized encoders; at $A_7$ the unchanged-objective
re-optimization remains compressed. A bootstrap caveat
travels with the joint rule: re-estimating $\hat\Sigma$ from
resampled honest sets (200 draws) leaves the gradient-family
rejection in 63\% of draws but flips the CE column between
all-pass and all-reject---at six-reference estimation the
calibration is reference-draw sensitive, and every conclusion
above is conditional on the primary instance (Supplementary
Table S12).



\paragraph*{The aggregate z-detector.}
\label{sec:audit-stat-implications}
The $z$-signal is specific to adversarially encoded residuals:
none of the three non-compliance constructions produces a
member-level signal, so a calibrated per-record statistic
separates deliberate encoding from sloppy unlearning on the
tested variants. The measured aggregate form---a one-sided
top-15 $z$-mean over $F$---detects 10 of 12 encoded finals
while firing on none of six non-compliance models (H0 control
1/2); a confirmatory run on MobileNetV2 finals detects 6 of 6. The one-sided variant was chosen after the frozen two-sided
gate failed (post-hoc); the MobileNetV2 run is the
confirmatory anchor. On a held-out honest evaluation (12
fresh ft models per architecture, excluded from all
selection) the frozen rule fires 2/12 on both: a held-out false-positive
rate of 17\% (exact 95\% CI $[2\%, 48\%]$)---a point estimate
above the 5\% design level with an interval that contains
it, so the data neither validate a 5\% operating point nor
establish a violation; the detector remains a pilot, and the
two-sided variant has no power
(1/12)---member-direction aggregation carries the signal, and
the honest tail structure is the binding constraint. The
cross-class spillover (\S\ref{sec:zreadout}) adds that
monitoring should extend beyond the deleted class
(Supplementary Table S13). An adaptive arm against the
detector is now measured: encodings optimized under the
composite constraint (pass $A_6$ and keep the top-15
z-mean below the detector threshold) satisfy both rules on
all three tested starts, but deliver no CI-positive gain
($+0.047$ to $+0.077$, every interval crossing zero;
Supplementary Table S29)---the detector's constraint
suppresses signal delivery even where feasibility is
achievable, placing it alongside $A_7$ as
compression-robust against the tested evasion tier.

\paragraph*{Other tested defenses.} The one-sided fm cap
closes the mean-objective wash on CIFAR-10 but leaks on
CIFAR-100 (\S\ref{sec:fmband}); multi-subsample re-evaluation
catches $2/6$ washes; retained-probe randomization did not
block the tested encodings, consistent with utility being
the binding component; and under trace-free retrain
references with a six-of-six threshold, mean-difference and
pairwise objectives stall at $0.52$--$0.64$ (single
configuration; a references-vs-threshold confound is
unseparated).

\section{Discussion}\label{sec:discussion}

\subsection{What This Paper Established}
The introduction asked whether a dishonest provider can pass
every audit statistic while encoding per-record membership
signals. The answer is yes, under measured conditions
(Tables~\ref{tab:encoding-main}--\ref{tab:zreadout}): nine
of nine label-blind gains, no flag from the audit's own
membership statistic, and per-record TPR $0.43$ among
passing finals under the reference-calibrated readout. The
tolerance has a two-part
account---two joint closures whose causal reading (reference
family vs.\ start ancestry) the design leaves confounded, and
within the admitted fine-tune family a delivery bought with
the utility allowance and invariant to the forget-side width
(\S\ref{sec:allowance}). Whether
the auditor can fight back is answered conditionally: the
aligned seventh statistic and the aggregate z-detector hold
through every tested evasion tier, while the joint rule and
threshold tightening are each evaded by re-optimized encoders
(Section~\ref{sec:eval}).

\textbf{Why delivered strength depends on the tested audit.}
The primary accepted finals reach the lower $m_1$ band edge,
showing that this component constrains the observed
optimization trajectories. Within the tested instance,
widening the utility-side bands increases mean delivered
gain, whereas varying the forget-side bands over the
reported near-baseline range changes it little. These
measurements identify a conditional trade-off, not an
optimality bound: other encoders may improve signal at the
same allowance. The retrain-referenced negative remains
confounded with start ancestry, and changes to the allowance
must also be evaluated against honest acceptance.

\subsection{Relation to Prior Work}
\textbf{Agreements.} Aggregate acceptance and per-record
erasure being different axes is consistent with disagreeing
benchmark metrics~\cite{dorna2025openunlearning} and
per-example leakage~\cite{hayes2024inexact}; the per-scale
recalibration requirement aligns
with~\cite{goel2022adversarial}.
\textbf{Tensions.} Our rollback-exclusion result is stronger
than the implicit assumption in some audit designs that lazy
unlearning is a threat worth characterizing: on our grid,
every rollback path violates all five metrics simultaneously.
Our selection-control finding---matched-budget honest
selection accounts for all but $0.000$--$0.078$ of margin
recovery at exact-retrain references---suggests some reported
recovery scores overstate attack capability for want of this
control. We name no invalidated prior result (the relevant
work operates under different verification interfaces), but
the methodological point stands: a recovery score without a
selection-at-matched-budget control is uninterpretable.

\textbf{Non-contradictions.} Zhang et
al.~\cite{zhang2024fragile} showed that dishonest providers
can bypass replay-based verification. Their attack surface
(replay/backdoor protocols) and ours (banded behavioral
statistics) are different computational objects; both can
coexist in a single deployment.

\subsection{Limitations and Scope}\label{sec:limitations}
Each limitation is paired with the conclusion whose scope it
affects.

\emph{Single principal deleted class.} The principal
experiment targets one CIFAR-10 class (cat); the class
probes (frog, dog, airplane; Supplementary Table S23)
measure how far its conclusions transfer---label-blind
replication on frog, partial on dog, no positive interval
on airplane, and no working-point reproduction on any other
class. Neither half of the mechanism orders the classes
(airplane has a larger overall-accuracy half-width than the
working family and delivers nothing), so class dependence
remains an
empirical boundary rather than a predicted one
(\S\ref{sec:allowance}).

\emph{Classifier scale only.} Results cover CIFAR-10/100 with
ResNet-18 and MobileNetV2; no LLM-scale audit is tested,
bounding \emph{every quantitative claim}; the mechanism
argument (tolerance must cover honest dispersion) is
audit-family-independent, but LLM-scale feasibility of the
encoding channel is unverified.

\emph{Provider-known calibration.} The main experiments
grant the provider the audit's calibration; under
independently estimated calibration, the five-draw evidence
shows 12/15 finals pass the true audit, 8/15 carry strictly
positive gains, and 5/15 satisfy both, with joint success in
3/5 draws (two further gains are borderline under the
declared rule; Supplementary Tables S22, S28). This bounds the
\emph{deployability claim} to draw-level heterogeneity
rather than a systematic barrier.

\emph{Ranking-level Tier-1 harm.} The deployable label-blind
harm is ranking-level (AUC gains; TPR $\le 0.10$ under the
fixed decoder, Section~\ref{sec:eval}). Tier-2 reference
access strengthens readability at the measured operating
points, but reliable low-FPR identification is not established
by either reading---bounding the attack's \emph{practical
severity}.

\emph{Reference-infrastructure handling.} The Tier-2 readout
makes a passing model per-record readable; for the auditor
(which knows the deleted set) it adds nothing, so the
consequence is custody: the reference ensemble is a readout
key for parties without membership knowledge.

\emph{Persistence boundary.} The per-record readout does not
survive one epoch of benign fine-tuning (median retention
$0.0$, Section~\ref{sec:eval}): the readability belongs to
the returned artifact, not to a persistent watermark.

\emph{Fixed weak MIA probe.} The audit's $m_5$ uses a
two-feature logistic probe; the stronger-probe
re-measurement (\S\ref{sec:mitigation}, Supplementary
Table S30) rejects the frozen finals while adaptive encoding
survives and the Tier-2 readout still delivers (2/3 starts,
TPR up to $0.20$).

\subsection{Concrete Directions}\label{sec:open}
\emph{(1)~Harm-statistic banding.} The aligned seventh
statistic compressed the attack by two-thirds (honest
acceptance 9/12 vs.\ 8/12, $n{=}12$); the next step is an
\emph{optimal} aligned statistic, whose binding direction
coincides with the encoding's contrast
direction---a convex-optimization question over the reference
ensemble's covariance.

\emph{(2)~Cross-request audit design.} The sequential-request
probe showed that sequential protection inherits the
reference-operator's tolerance; a cross-request audit must
either accept this inheritance or budget fresh calibration
per request---a quantifiable trade-off.

\emph{(3)~LLM-scale probe.} The protocol transfers as an
experimental design: (i)~train a reference ensemble on
a small corpus-level deletion; (ii)~measure forget-side
dispersion relative to full-data bands; (iii)~where the
inflation factor is above $2\times$, treat the audit as worth
investigating---a screening threshold for where to spend a
costly probe, not a trigger that expects encoding to succeed
(at classifier scale $k{=}100$ shows $2.5\times$ inflation
with no deliverable signal; Supplementary Table S36).


\section{Conclusion and Future Work}\label{sec:conclusion}
This paper proposed ECLIPSE, to our knowledge the first
construction of deliberate per-record membership encoding
that passes a live, conformally calibrated multi-metric
deletion audit. FIDE constructs the encoded models and
boundary pinning characterizes the accepted endpoints:
label-blind scoring establishes the ranking gains (nine of
nine starts), and the reference-equipped CAL-Z protocol the
working points (per-record TPR $0.43$ among passing finals
at Tier 2).

The mechanism account is two-part and bounded (Section~VI):
admissibility closes the retrain family, the utility allowance
governs delivery within the admitted family, and neither
orders the class or scale axes. The audit's defenses
split by instrument: the aligned seventh statistic and the
aggregate z-detector hold through every tested evasion tier,
while the joint rule and threshold tightening are each evaded
by re-optimized encoders---an arms race that moves with the
rule rather than ending it.

The measurement discipline introduced here---matched-budget
selection controls, logit-scale controls,
dispersion-calibrated bands, and preregistered adaptive
counter-arms---is itself a contribution: it is what
distinguishes a delivered per-record signal from
margin-recovery readings that honest processing already
produces. For practitioners the lesson is direct: the
certificate should state its coverage scope, its
calibration-knowledge assumptions, and what it excludes.
Acceptance is not erasure: what an audit aggregates
determines what it can be deceived about. Future work
follows the directions of Section~\ref{sec:open}: an optimal
aligned harm statistic, cross-request audit design, and the
LLM-scale probe.


\bibliographystyle{IEEEtran}

\IEEEtriggeratref{16}
\IEEEtriggercmd{\enlargethispage{7\baselineskip}}

\bibliography{refs}

\end{document}
